\documentclass[12pt]{article}

\usepackage{amsmath, amssymb, amsfonts}
\usepackage{geometry}
\usepackage{booktabs}
\usepackage{setspace}
\usepackage{graphicx}
\usepackage{hyperref}
\usepackage{longtable}
\usepackage{array}
\usepackage{enumitem}
\usepackage{float}
\usepackage{multirow}
\usepackage{caption}
\usepackage{pdflscape}

\geometry{margin=1in}
\onehalfspacing

\hypersetup{
    colorlinks=true,
    linkcolor=blue,
    citecolor=blue,
    urlcolor=blue
}

\begin{document}

% ============================================================
% TITLE PAGE
% ============================================================
\begin{titlepage}
\centering
\vspace*{2cm}

{\LARGE \textbf{Nonparametric Statistical Analysis of\\[0.3cm] Statewide Planning and Research Cooperative System 2022 Hospital Discharge Data:\\[0.3cm] Kings and Queens Counties, New York}}

\vspace{1.0cm}

{\large Combined Midterm and Final Project Report}

\vspace{1.5cm}

{\Large Vinay Siddi}

\vspace{1cm}

{\large AMAT 581 -- Nonparametric Statistics}

\vspace{0.5cm}

{\large Spring 2026}

\vspace{1cm}

{\large Department of Data Science\\
University at Albany, SUNY}

\vspace{2cm}

{\large \today}

\vfill
\end{titlepage}

\newpage
\tableofcontents
\newpage

% ============================================================
% ABSTRACT (merged)
% ============================================================
\section*{Abstract}
\addcontentsline{toc}{section}{Abstract}

This project applies ten distribution-free (nonparametric) statistical
procedures to the New York State SPARCS (Statewide Planning and Research
Cooperative System) 2022 hospital inpatient discharge data set, restricted
to Kings County (Brooklyn) and Queens County. The data set contains
368{,}495 de-identified discharge records on 33 original Statewide Planning and Research Cooperative System variables
(plus 4 derived numeric columns created during cleaning, for 37 total
columns in the analysis data frame) and is approximately 160~MB as a Comma-Separated Values
file on disk (136.9~MB as an R in-memory data frame). Four quantitative
outcomes are central to the analysis: Length of Stay (days), Total Charges
(US Dollars), Total Costs (US Dollars), and Birth Weight (grams).

The tests performed are (1) the \textbf{Runs Test} for randomness,
(2) the \textbf{Sign Test}, (3) the \textbf{Wilcoxon Signed-Rank Test},
(4) the \textbf{Wilcoxon Rank Sum Test (Mann--Whitney~U)}, (5) the
two-sample \textbf{Kolmogorov--Smirnov Test}, (6) the
\textbf{Kruskal--Wallis Test} for $k$ independent groups, (7) the
\textbf{Friedman Test} with facility as block and All Patient Refined Severity of Illness
as treatment, (8) \textbf{Kendall's Tau}, (9) \textbf{Spearman's Rank
Correlation}, and (10) the \textbf{Permutation Test} (Monte Carlo,
$B=5{,}000$). As an \textbf{innovative contribution}, the study performs
a simultaneous application of all these tests to a single large Statewide Planning and Research Cooperative System
slice (368{,}495 records), combines $p$-value significance testing with
visualization of the empirical permutation null distribution, and extends
the analysis to interactions between the Kings/Queens county label and
admission severity, which no prior work on the Statewide Planning and Research Cooperative System has attempted at this
scale.

All principal tests reject the null hypothesis across essentially every
meaningful comparison: the raw record ordering is non-random (Runs Test);
the median and distribution of Length of Stay differ significantly
between Kings and Queens counties (Sign Test $p<10^{-16}$;
Wilcoxon Signed-Rank and Mann--Whitney~U both $p<10^{-16}$); the entire
cumulative distributions of Length of Stay, Charges, Costs, and Birth Weight differ
between the two counties (Kolmogorov-Smirnov $p<10^{-16}$); every numeric variable
differs across Age Group, Severity, Admission Type, and Payer
(Kruskal--Wallis $p<10^{-16}$); and Length of Stay, Charges, and Costs differ across
severity levels even after blocking by facility (Friedman $p<0.005$).
The strongest monotone association is between Charges and Costs
(Spearman $\rho=0.86$); Length of Stay--Charges is also strong ($\rho=0.73$).
The permutation test on Total Charges yields a significant median
difference in Charges between counties ($p=0.011$), and the
large-sample Kolmogorov-Smirnov and Kruskal--Wallis tests on Length
of Stay reject resoundingly. All analyses were implemented in base R;
no external packages were used for the test statistics themselves.

% ============================================================
% 1. INTRODUCTION (merged)
% ============================================================
\section{Introduction}

New York State's Statewide Planning and Research Cooperative System database is one of the most granular all-payer
hospital discharge data sets in the United States, covering essentially
every inpatient admission in the state. Within New York City, Kings
County (Brooklyn) and Queens County together account for the largest
share of non-Manhattan admissions, serving substantially different
demographic populations. Length of Stay (Length of Stay), Total Charges, and Total
Costs distributions on hospital data are well known to be heavily
right-skewed with long tails, which makes the usual parametric $t$-tests
and ANOVAs inappropriate. Nonparametric (distribution-free) methods are
the standard tool for comparing such data.

This project applies ten nonparametric tests (the Runs Test, Sign Test,
Wilcoxon Signed-Rank Test, Mann--Whitney~U Test, Kolmogorov--Smirnov
Test, Kruskal--Wallis Test, Friedman Test, Kendall's Tau, Spearman's
Rank Correlation, and Permutation Test) to quantify:
(i) whether the raw ordering of the Statewide Planning and Research Cooperative System records is random; (ii)
whether the median and distribution of Length of Stay/Charges/Costs/Birth~Weight
differ between the two counties and between Emergency and Non-Emergency
admissions; (iii) whether these outcomes differ across sub-groups (age,
severity, admission type, payer); (iv) whether severity drives
utilization even after controlling for facility; (v) whether the four
quantitative outcomes are monotonically associated; and (vi) whether the
county-label differences are attributable to chance under an
exchangeability null.

% ============================================================
% 2. LITERATURE REVIEW (merged)
% ============================================================
\section{Literature Review}

Nonparametric statistical analysis of hospital discharge data has grown
a lot as a research area since 2010. Three main reasons drove this
growth: first, researchers realized that Length of Stay, hospital cost,
and hospital charge values are almost always heavily skewed to the
right (meaning most values are small but a few are extremely large),
which breaks the assumptions that classical tests like the t-test and
ANOVA rely on. Second, several large public hospital databases became
available, including the Statewide Planning and Research Cooperative
System in New York, the Healthcare Cost and Utilization Project
Nationwide Inpatient Sample at the national level, and the National
Hospital Ambulatory Medical Care Survey for outpatient visits. Third,
more methodology papers started appearing that used simulations to
compare how often different tests give wrong answers (false positives
and false negatives) on realistic hospital data. The papers below
motivate the present project and are presented in chronological order
of publication.

\textbf{Swed and Eisenhart (1943)} published the original tables of
critical values for the Runs Test in the \textit{Annals of
Mathematical Statistics} (Volume 14, Number 1, pages 66 to 87). Their
methods combined exact combinatorial enumeration of the sampling
distribution of the number of runs with the large-sample
normal-approximation framework, producing the critical-value tables
that are still used today. The normal-approximation formula I use in
the Runs Test section of this project comes directly from their
theoretical results.

\textbf{Efron and Tibshirani (1993)} wrote \textit{An Introduction to
the Bootstrap}, published by Chapman and Hall / CRC. They introduced
the bootstrap, which is a resampling method for estimating
uncertainty. Their book covers the nonparametric bootstrap, the
parametric bootstrap, the jackknife, and bootstrap confidence
intervals (percentile, bias-corrected, and BCa methods), positioning
all of these against classical parametric standard-error formulas.
Although bootstrap confidence intervals are only mentioned as a
future extension in my Conclusions, the permutation test that I do
run in this project is in the same family of resampling-based methods.

\textbf{Ernst (2004)} published a tutorial in \textit{Statistical
Science} (Volume 19, Number 4, pages 676 to 685) explaining
permutation methods as an alternative to classical parametric tests
such as the Student t-test, the F-test, and Pearson correlation
testing. He argued that permutation tests are exact when the null
hypothesis of exchangeable labels holds, do not require any assumption
about the shape of the distribution, and are especially good for
modern large datasets. This tutorial is the direct methodological
foundation for the Permutation Test in my project.

\textbf{Fagerland and Sandvik (2009)} published an examination of the
Wilcoxon-Mann-Whitney test in \textit{Statistics in Medicine} (Volume
28, pages 1487 to 1497). Their methods combined Monte Carlo simulation
studies of the Wilcoxon-Mann-Whitney rank-sum test with comparisons to
the Student t-test and the Welch t-test under unequal variances and
unequal distribution shapes. They showed that even moderate
differences in group variance or distribution shape can cause the
Wilcoxon-Mann-Whitney test to reject more often than it should, but
the test is still useful when the alternative is a straight shift
between the two groups. This warning is relevant to my Kings versus
Queens comparison, where the within-county variance of charges is
very different between the two groups. This is exactly why I included
the Kolmogorov-Smirnov test in this project as a check on the whole
distribution shape.

\textbf{Qualls, Pallin, and Schuur (2010)} published a study in
\textit{Academic Emergency Medicine} (Volume 17, Number 10, pages 1113
to 1121). They looked at every Emergency Department study in five
major Emergency Medicine journals between 2004 and 2007 that reported
Length of Stay, and they catalogued the statistical methods those
studies used: Student's t-test, one-way and two-way ANOVA, linear
regression, and Pearson correlation as parametric methods, alongside
the Wilcoxon rank-sum, Mann-Whitney U, Kruskal-Wallis, and chi-squared
tests as nonparametric alternatives. They found that 43 percent of
these studies used classical parametric methods even though Length of
Stay data is clearly not normally distributed. Using 2006 National
Hospital Ambulatory Medical Care Survey data, they showed that this
mistake causes false negatives (missing real differences) in simple
two-group comparisons and weakens the explanatory power of more
complex models. Their paper is the most cited justification for using
Wilcoxon, Mann-Whitney, and Kruskal-Wallis tests on hospital Length
of Stay data, and it is a direct motivation for the present project.

\textbf{Croux and Dehon (2010)} published a study in
\textit{Statistical Methods and Applications} (Volume 19, pages 497
to 515). Their methods combined influence-function analysis of
Spearman's Rho and Kendall's Tau with theoretical derivations under
the bivariate-normal distributional model, plus comparisons with the
classical Pearson product-moment correlation coefficient. They proved
that Kendall's Tau is more robust to outliers than Spearman's Rho
because its influence function is bounded. They also showed that,
under a two-variable normal distribution, Tau is approximately two
over pi times the arc sine of Rho. This means Tau will always be
smaller in size than Rho for the same pair of variables. This is
exactly the pattern I see across all six numeric pairs in my Kendall's
Tau and Spearman sections of the project.

\textbf{Anhoj and Olesen (2014)} published a simulation study in
\textit{PLOS ONE} (Volume 9, Number 11, article e113825). Their
methods combined Monte Carlo simulation of run-chart rules with
Shewhart control-chart theory and runs-test theory to evaluate how
well these tools detect non-random variation in healthcare processes.
They found that rules based on runs analysis have a detection rate of
around 95 percent for process shifts, while the false-alarm rate stays
near 5 percent regardless of how many observations are in the sample.
Their work shows that runs-based methods are reliable tools for
healthcare quality improvement.

\textbf{Moore, White, Washington, Coenen, and Elixhauser (2017)}
published a study in \textit{Medical Care} (Volume 55, Number 7,
pages 698 to 705). They used the Healthcare Cost and Utilization
Project Nationwide Inpatient Sample, which is the federal counterpart
to the Statewide Planning and Research Cooperative System, to
identify demographic and clinical predictors of readmission risk and
in-hospital death. Their methods combined parametric multivariable
logistic regression and the AHRQ Elixhauser comorbidity index with
nonparametric Wilcoxon rank-sum and chi-squared tests for group
comparisons. They showed that administrative hospital cost
distributions differ significantly by demographic group and by
admission type. This matches the payer-group and age-group effects
that my project measures on the Kings and Queens county data.

\textbf{Chazard, Ficheur, Beuscart, and Preda (2017)} published a
simulation study in \textit{Value in Health} (Volume 20, Number 7,
pages 992 to 998). Their methods compared twelve different statistical
tests on simulated Length-of-Stay-like data drawn from a realistic
right-skewed distribution with many repeated values. The twelve tests
included the Student t-test, the Welch t-test, ANOVA, the
Wilcoxon-Mann-Whitney rank-sum test, the Kruskal-Wallis test, the
median test, log-rank tests, bootstrap resampling, and permutation
methods. They found that the Student t-test gives false positive rates
above 10 percent when the group variances are unequal (much higher
than the intended 5 percent), while the Wilcoxon/Mann-Whitney test
keeps its false positive rate at the correct level and still has good
power. They also showed that bootstrap and permutation methods hold up
well when the distributions are very skewed. Their simulation results
directly support my choice of the Kolmogorov-Smirnov, Kruskal-Wallis,
and Permutation tests in this project.

\textbf{Ning, Harkness, and Gao (2017)} published a study in the
\textit{American Medical Informatics Association Annual Symposium
Proceedings} (PubMed Central ID PMC5543354). Their methods combined
spatio-temporal analysis using Moran's I statistic with classical
geostatistical clustering techniques applied to Statewide Planning
and Research Cooperative System data from 2003 to 2014 at the
zip-code level. Their work shows how rich the Statewide Planning and
Research Cooperative System database is for geographic comparisons of
hospital use. My project narrows that same kind of geographic
comparison down to the county level.

\textbf{Bujang and Sapri (2018)} published a study in the
\textit{Malaysian Journal of Medical Sciences} (Volume 25, Number 4,
pages 146 to 151). Their methods applied the Runs Test (using the
normal-approximation Z-statistic) to clinical hemoglobin A1c data
collected by consecutive sampling, with descriptive parametric
summaries (sample mean, standard deviation) reported alongside the
nonparametric randomness test. They recommended that randomness
testing should be done in every clinical pilot study. They gave a
clear example showing that data collected one patient after another
usually does not come out in random order. This matches my own Runs
Test results: the Statewide Planning and Research Cooperative System
records are non-random because they are stored in facility and
admission-date order.

\textbf{Dieleman, Cao, Chapin, and colleagues (2020)} published a
study in the \textit{Journal of the American Medical Association}
(Volume 323, Number 9, pages 863 to 884). They measured United States
healthcare spending from 1996 to 2016, broken down by payer type and
by health condition. Their methods combined parametric decomposition
analysis (a regression-based attribution of spending to price,
quantity, and service-mix factors) with nonparametric bootstrap
confidence intervals to quantify spending trends. They found big
differences in hospital charges and costs across U.S. regions. Their
findings set the baseline expectation that county-level differences
in charges should be detectable in the Statewide Planning and
Research Cooperative System data, which is the main hypothesis that
my project tests.

\textbf{Nguyen, Ho, Nguyen, and Van (2021)} published a study in the
\textit{VNUHCM Journal of Engineering and Technology} (Volume 3,
Special Issue 3, pages SI97 to SI106). They analyzed 1,189 Vietnamese
patient records using a mix of parametric and nonparametric methods,
including the Student t-test, one-way ANOVA, Pearson correlation, and
multiple linear regression on the parametric side, alongside the
Kruskal-Wallis Test and the Mann-Whitney U Test on the nonparametric
side. They found that age, province, profession, admission quarter,
and type of disease are the main factors that affect Length of Stay.
Their findings about age and severity driving Length of Stay
distributions are reproduced on the much larger Statewide Planning
and Research Cooperative System dataset in the present report.

\textbf{Stone, Zwiggelaar, Jones, and Mac Parthalain (2022)} published
a systematic review in \textit{PLOS Digital Health} (Volume 1, Number
4, article e0000017). They reviewed 93 studies from 1970 to 2019 that
tried to predict hospital Length of Stay and catalogued the methods
those studies used: parametric statistical methods (linear regression,
generalized linear models, Cox proportional-hazards regression, log-normal
and gamma regression), nonparametric statistical methods (Kaplan-Meier
survival curves, log-rank tests), and machine-learning / data-mining
methods (decision trees, random forests, support vector machines,
neural networks). Their most striking numerical finding is that
Length of Stay for the same diagnosis can range anywhere from 2 days
to more than 50 days. This huge variation is what makes the
Kruskal-Wallis test useful for my project, where I compare Length of
Stay across severity levels and admission types.

\textbf{Li, Couper, Hunt, and Mian (2022)} published a study in the
\textit{Journal of the Royal Statistical Society: Series C (Applied
Statistics)} (Volume 71, Number 5, pages 1667 to 1693). Their methods
combined a parametric Weibull-type hazard function with a
nonparametric baseline curve to build a semi-parametric time-to-event
model for Length of Stay that simultaneously models the chance of
discharge and the chance of death; they also reported maximum
likelihood estimation and likelihood-ratio tests. They tested their
method on stroke clinical trial data and on COVID-19 patient data.
Their hybrid approach is listed in my Conclusions section as a
possible future extension of the present analysis.

\textbf{Heslin and colleagues (2024)} published a study in
\textit{Cureus} (Volume 16, Number 4, article e58404). Their methods
combined chi-squared tests of independence with one-way analysis of
variance (ANOVA) and post-hoc Tukey HSD comparisons to evaluate
admission and discharge patterns by day of the week using Statewide
Planning and Research Cooperative System data covering 1,994,180
adult hospital admissions from January to December 2015. They found
that elective surgical admissions explain almost all of the weekly
variation in hospital capacity. Hospitals ran 21 percent above
average volume on Mondays and 32 percent below average on Fridays,
with a p-value less than 0.01. Their study showed that the Statewide
Planning and Research Cooperative System database is well suited for
large-scale nonparametric analyses, which directly supports the
approach I take in this project.

Taken together, the papers above make three main points that frame
my work:

\begin{enumerate}[nosep]
  \item Hospital Length of Stay, hospital charges, and hospital costs
    are all heavily right-skewed, so they must be analyzed with
    nonparametric methods to avoid inflated false-positive or
    false-negative rates (Qualls, Pallin, and Schuur, 2010;
    Chazard and colleagues, 2017; Fagerland and Sandvik, 2009).
  \item The Statewide Planning and Research Cooperative System is a
    well-established, high-quality administrative data source for
    these analyses. Earlier Statewide Planning and Research Cooperative System studies have documented big
    variation across day of the week, geography, and admission type
    (Heslin and colleagues, 2024; Ning, Harkness, and Gao, 2017).
  \item The theoretical foundation for the ten nonparametric tests
    I use is laid out in the specific methodology papers for each
    test (Swed and Eisenhart, 1943; Croux and Dehon, 2010; Ernst,
    2004; Efron and Tibshirani, 1993).
\end{enumerate}

My project contributes to the literature by applying \emph{ten}
nonparametric tests at the same time to a single large
(368,495-record) publicly available Statewide Planning and Research
Cooperative System dataset. The tests cover randomness (Runs Test),
center of location (Sign Test, Wilcoxon Signed-Rank Test, Mann-Whitney
U Test), full distribution equality (Kolmogorov-Smirnov Test),
group-difference (Kruskal-Wallis Test, Friedman Test), rank
association (Kendall's Tau, Spearman's Rank Correlation), and label
exchangeability (Permutation Test). To the best of my knowledge, this
is the first comprehensive distribution-free comparison of hospital
use between Kings County and Queens County on the Statewide Planning
and Research Cooperative System 2022 dataset.

% ============================================================
% 3. DATA PRESENTATION AND EXPLORATORY ANALYSIS
% ============================================================
\section{Data Presentation and Exploratory Analysis}

\subsection{Dataset Overview}

The primary data set is \texttt{SPARCS\_2022\_Kings\_Queens.csv}, a
filtered subset of the publicly available New York State SPARCS
Inpatient Discharges (De-identified) 2022 data, restricted to the two
counties of interest.

\begin{itemize}[nosep]
    \item \textbf{Number of Observations:} 368{,}495 rows.
    \item \textbf{Number of Variables:} 33 original Statewide Planning and Research Cooperative System columns, plus
      4 derived numeric columns (\texttt{Length of Stay}, \texttt{Charges},
      \texttt{Costs}, \texttt{BirthWt}) created during cleaning, for 37
      total columns in the analysis data frame.
    \item \textbf{File Size:} approximately 160~MB as a Comma-Separated Values file on disk
      (136.9~MB as an R in-memory data frame via \texttt{object.size()}).
    \item \textbf{Time Period:} January 1, 2022 -- December 31, 2022.
    \item \textbf{Geographic Coverage:} Kings and Queens Counties, New York City.
    \item \textbf{Unit of Observation:} individual hospital discharge record.
    \item \textbf{Source:} New York State Department of Health,
      Statewide Planning and Research Cooperative System (SPARCS),
      \url{https://health.data.ny.gov/}.
\end{itemize}

\subsection{Variable Descriptions}

The 33 Statewide Planning and Research Cooperative System variables can be classified by measurement scale (rather than by R storage type) into four truly quantitative (ratio-scale) variables — Length of Stay, Total Charges, Total Costs, and Birth Weight; one ordinal variable — All Patient Refined Severity of Illness Code (1--4); and 28 categorical (nominal) variables. Note that several of the categorical variables (Operating Certificate Number, Permanent Facility ID, All Patient Refined Diagnosis-Related Group Code, and All Patient Refined Major Diagnostic Category Code) are stored as integers in the raw file, but their numeric values are arbitrary identifiers or classification labels — arithmetic on them (e.g.\ a mean) is not meaningful, so they are classified here as nominal. Discharge Year is stored as an integer but is constant (2022) for every record. Table~\ref{tab:variables} provides a detailed listing.

\begin{longtable}{p{0.22\textwidth} p{0.16\textwidth} p{0.32\textwidth} p{0.22\textwidth}}
\caption{Variables in the Statewide Planning and Research Cooperative System 2022 Kings and Queens Dataset}
\label{tab:variables} \\
\toprule
\textbf{Variable Name} & \textbf{Type} & \textbf{Meaning of Variables} & \textbf{Unit / Description} \\
\midrule
\endfirsthead
\multicolumn{4}{c}{\textit{(continued)}} \\
\toprule
\textbf{Variable Name} & \textbf{Type} & \textbf{Meaning of Variables} & \textbf{Unit / Description} \\
\midrule
\endhead
Hospital Service Area & Qualitative (Categorical) & Geographic service area name reported by the facility, used to group hospitals serving the same regional population. & Service area name (text) \\
Hospital County & Qualitative (Categorical) & County in which the hospital is located; the central grouping variable in this study. & Kings or Queens (text) \\
Operating Certificate Number & Qualitative (Nominal) & State-issued certificate identifier that authorizes the facility to operate. & Facility certificate ID (integer label) \\
Permanent Facility ID & Qualitative (Nominal) & Permanent unique identifier assigned to each hospital facility by the state. & Unique facility identifier (integer label) \\
Facility Name & Qualitative (Categorical) & Name of the discharging hospital. & Hospital name (text) \\
Age Group & Qualitative (Categorical) & Patient age binned into broad cohorts to allow group comparisons while protecting privacy. & 0--17, 18--29, 30--49, 50--69, 70+ (text) \\
Zip Code (3 digits) & Qualitative (Categorical) & First three digits of the patient's home ZIP code, used as a coarse geographic indicator. & First 3 digits of patient zip (text) \\
Gender & Qualitative (Categorical) & Patient gender as recorded at admission. & F, M, U (text) \\
Race & Qualitative (Categorical) & Patient self-reported race. & Patient race category (text) \\
Ethnicity & Qualitative (Categorical) & Patient self-reported Hispanic / Latino ethnicity. & Patient ethnicity (text) \\
Length of Stay & Quantitative (Continuous) & Number of days from admission to discharge; one of the four primary outcome variables. & Days hospitalized (days) \\
Type of Admission & Qualitative (Categorical) & How the patient was admitted (emergency vs.\ planned vs.\ urgent, etc.). & Emergency, Elective, Urgent, etc.\ (text) \\
Patient Disposition & Qualitative (Categorical) & Where the patient went after discharge (home, skilled nursing, expired, etc.). & Discharge status (text) \\
Discharge Year & Qualitative (Constant) & Calendar year of discharge; constant at 2022 for the slice analyzed in this project. & 2022 for every record (integer, no variation) \\
CCSR Diagnosis Code & Qualitative (Categorical) & Clinical Classifications Software Refined diagnosis category code. & Clinical classification code (text) \\
CCSR Diagnosis Description & Qualitative (Categorical) & Plain-text description matching the CCSR diagnosis code. & Diagnosis description (text) \\
CCSR Procedure Code & Qualitative (Categorical) & CCSR procedure category code for any procedures performed. & Procedure classification code (text) \\
CCSR Procedure Description & Qualitative (Categorical) & Plain-text description matching the CCSR procedure code. & Procedure description (text) \\
All Patient Refined Diagnosis-Related Group Code & Qualitative (Nominal) & Numeric DRG label that groups patients by clinical similarity for billing and analysis. & Diagnosis-Related Group classification label (integer) \\
All Patient Refined Diagnosis-Related Group Description & Qualitative (Categorical) & Plain-text description matching the DRG code. & Diagnosis-Related Group description (text) \\
All Patient Refined Major Diagnostic Category Code & Qualitative (Nominal) & High-level grouping of DRGs into broad body-system categories. & Major Diagnostic Category label (integer) \\
All Patient Refined Major Diagnostic Category Description & Qualitative (Categorical) & Plain-text description matching the MDC code. & Major Diagnostic Category description (text) \\
All Patient Refined Severity of Illness Code & Quantitative (Ordinal) & Ordinal severity rating from Minor (1) to Extreme (4); used as the treatment factor in the Friedman test. & 1 = Minor, 2 = Moderate, 3 = Major, 4 = Extreme \\
All Patient Refined Severity of Illness Desc.\ & Qualitative (Categorical) & Text version of the severity-of-illness rating. & Minor, Moderate, Major, Extreme (text) \\
All Patient Refined Risk of Mortality & Qualitative (Categorical) & Ordinal indicator of how likely the patient is to die during the admission. & Risk level description (text) \\
All Patient Refined Medical Surgical Desc.\ & Qualitative (Categorical) & Whether the encounter is classified as a medical or a surgical case. & Medical or Surgical (text) \\
Payment Typology 1 & Qualitative (Categorical) & Primary expected payer for the admission (Medicare, Medicaid, Private, etc.). & Primary payer (text) \\
Payment Typology 2 & Qualitative (Categorical) & Secondary expected payer when more than one payer is responsible. & Secondary payer (text) \\
Payment Typology 3 & Qualitative (Categorical) & Tertiary expected payer when applicable. & Tertiary payer (text) \\
Birth Weight & Quantitative (Continuous) & Weight at birth in grams, populated only for newborn discharges. & Weight in grams (newborns only) \\
Emergency Dept.\ Indicator & Qualitative (Categorical) & Yes/No flag for whether the admission came through the Emergency Department. & Y or N (text) \\
Total Charges & Quantitative (Continuous) & Total amount billed for the admission; one of the four primary outcome variables. & Hospital charges (US Dollars) \\
Total Costs & Quantitative (Continuous) & Estimated actual cost incurred by the hospital for the admission; one of the four primary outcome variables. & Hospital costs (US Dollars) \\
\bottomrule
\end{longtable}

\subsection{Levels of the Principal Qualitative Variables}

Table~\ref{tab:levels} records the distinct levels observed in the key
categorical variables, which drive the Kruskal--Wallis
stratifications.

\begin{table}[H]
\centering
\caption{Levels of Key Qualitative Variables (Statewide Planning and Research Cooperative System 2022 K\&Q)}
\label{tab:levels}
\renewcommand{\arraystretch}{1.15}
\begin{tabular}{p{0.22\textwidth} c p{0.66\textwidth}}
\toprule
\textbf{Variable} & \textbf{\# Levels} & \textbf{Levels (with descriptions)} \\
\midrule
Hospital County & 2 &
\textbf{Kings:} Brooklyn-based hospitals serving Kings County residents.
\newline
\textbf{Queens:} Queens-based hospitals serving Queens County residents. \\
\midrule
Age Group & 5 &
\textbf{0 to 17:} Pediatric patients (children and adolescents).
\newline
\textbf{18 to 29:} Young adults; typically the healthiest cohort.
\newline
\textbf{30 to 49:} Working-age adults.
\newline
\textbf{50 to 69:} Middle-aged to early-elderly adults.
\newline
\textbf{70 or Older:} Elderly patients; typically the highest-utilization cohort. \\
\midrule
Gender & 3 &
\textbf{Female:} Recorded as female at admission.
\newline
\textbf{Male:} Recorded as male at admission.
\newline
\textbf{Unidentified:} Gender not recorded or not classified at admission. \\
\midrule
All Patient Refined Severity of Illness Code & 4 &
\textbf{1 -- Minor:} Lowest severity; routine cases with no significant complications.
\newline
\textbf{2 -- Moderate:} Moderate severity; some complications or comorbidities.
\newline
\textbf{3 -- Major:} High severity; major complications expected to extend the stay.
\newline
\textbf{4 -- Extreme:} Highest severity; life-threatening illness with high resource use. \\
\midrule
All Patient Refined Risk of Mortality & 4 &
\textbf{Minor:} Lowest predicted likelihood of in-hospital death.
\newline
\textbf{Moderate:} Some elevation in mortality risk above baseline.
\newline
\textbf{Major:} Substantially elevated mortality risk.
\newline
\textbf{Extreme:} Very high predicted likelihood of in-hospital death. \\
\midrule
Type of Admission & 5 &
\textbf{Elective:} Planned admission, scheduled in advance.
\newline
\textbf{Emergency:} Unscheduled admission requiring immediate care.
\newline
\textbf{Newborn:} Admission of a newborn at or shortly after delivery.
\newline
\textbf{Urgent:} Unscheduled but not immediately life-threatening; care needed within 24--48 hours.
\newline
\textbf{Trauma:} Admission resulting from acute physical injury (accident, assault, etc.). \\
\midrule
Emergency Department Indicator & 2 &
\textbf{Yes:} The patient was admitted through the Emergency Department.
\newline
\textbf{No:} The patient was admitted through a non-Emergency-Department route (direct, planned, transfer, etc.). \\
\midrule
Payment Typology 1 & $\approx$\,9 &
\textbf{Medicare:} Federal health insurance program for those aged 65+ and certain younger people with disabilities.
\newline
\textbf{Medicaid:} Joint federal-state program covering low-income individuals and families.
\newline
\textbf{Private Insurance:} Commercial insurance through an employer or purchased individually.
\newline
\textbf{Self-Pay:} The patient is personally responsible for the bill, with no insurance coverage.
\newline
\textbf{Other:} Remaining categories such as Workers' Compensation, Federal/State/Local government plans, Department of Corrections, miscellaneous, and unknown payers. \\
\midrule
All Patient Refined Medical Surgical Description & 2 &
\textbf{Medical:} The encounter is classified as a non-surgical (medical) case.
\newline
\textbf{Surgical:} The encounter involved a surgical procedure during the admission. \\
\bottomrule
\end{tabular}
\end{table}

\subsection{Summary Statistics of Key Numeric Variables}

\begin{table}[H]
\centering
\caption{Summary Statistics of Key Numeric Variables (Statewide Planning and Research Cooperative System 2022)}
\label{tab:summary}
\begin{tabular}{lrrrrrr}
\toprule
\textbf{Variable} & \textbf{Min} & \textbf{Q1} & \textbf{Median} & \textbf{Mean} & \textbf{Q3} & \textbf{Max} \\
\midrule
Length of Stay (days)  & 1 & 2 & 3 & 5.76 & 6 & 119 \\
Total Charges (US Dollars)    & 7 & 21{,}544 & 40{,}497 & 71{,}907 & 77{,}905 & 12{,}952{,}195 \\
Total Costs (US Dollars)      & 2 & 6{,}746 & 12{,}543 & 23{,}273 & 24{,}821 & 5{,}602{,}076 \\
Birth Weight (grams)   & 400 & 2{,}900 & 3{,}200 & 3{,}159 & 3{,}500 & 6{,}100 \\
\bottomrule
\end{tabular}
\end{table}

\textit{Note: Length of Stay has 332 missing values. Birth Weight has
325{,}132 missing values (only newborn discharges). Total Charges and
Total Costs have no missing values.}

\subsection{Exploratory Graphs}

The accompanying R~Markdown file produces the following reproducible
graphs. Each figure is shown below alongside a brief interpretation;
the underlying R code appears in the companion \texttt{.Rmd} file.

\begin{figure}[H]
\centering
\includegraphics[width=0.85\textwidth]{figures/fig1_los_county.png}
\caption{Log-scale boxplot of Length of Stay by Hospital County. Both
counties show very heavy right skew, with Kings extending slightly
higher than Queens.}
\label{fig:los-county}
\end{figure}

\textbf{What this plot shows.} This boxplot compares how long patients
stay in the hospital between Kings County and Queens County. The thick
black line in each box is the median (the typical patient's length of
stay), and the box itself covers the middle 50\% of patients. The
small dots above each box are unusually long stays, and there are many
of them in both counties because a small number of very sick patients
stay in the hospital for weeks. Most patients in both counties go home
in 1 to 5 days, with Kings County stays lasting slightly longer on
average than Queens County stays.

\begin{figure}[H]
\centering
\includegraphics[width=0.85\textwidth]{figures/fig2_charges_county.png}
\caption{Log-scale boxplot of Total Charges by Hospital County. The
median Charges in Kings County are visibly higher than in Queens
County, consistent with the Mann--Whitney and Kolmogorov--Smirnov
results.}
\label{fig:charges-county}
\end{figure}

\textbf{What this plot shows.} This boxplot compares the total
hospital charges (the dollar amount billed) between Kings and Queens
counties. The y-axis uses a log scale because charges range from
under \$100 to over \$10 million, so a regular scale would make small
charges invisible. The typical patient is billed somewhere around
\$40,000 in both counties, but Kings County's median is noticeably
higher than Queens County's. Many extreme outliers extending toward
\$10 million represent the very expensive cases like long ICU stays
or complex surgeries.

\begin{figure}[H]
\centering
\includegraphics[width=0.85\textwidth]{figures/fig3_costs_county.png}
\caption{Log-scale boxplot of Total Costs by Hospital County. The same
upward shift is visible in Costs as in Charges.}
\label{fig:costs-county}
\end{figure}

\textbf{What this plot shows.} This boxplot compares total hospital
costs (the actual cost to the hospital, not the amount billed) between
the two counties. Costs are typically much lower than charges because
charges are inflated for negotiation with insurance companies. The
pattern is the same as for charges: Kings County tends to have
slightly higher costs than Queens County, and a small number of
extreme cases push costs into the millions. The strong similarity
between the charge and cost patterns suggests that hospitals in both
counties use similar markup ratios.

\begin{figure}[H]
\centering
\includegraphics[width=0.95\textwidth]{figures/fig_los_histograms.png}
\caption{Histograms of Length of Stay (capped at 30~days) for Kings
and Queens. More than 50\% of discharges have Length of
Stay~$\le$~3~days in both counties, confirming extreme right-skew.}
\label{fig:los-histograms}
\end{figure}

\textbf{What this plot shows.} These two histograms show how often
each length of stay occurs in each county, capped at 30 days so the
shape stays readable. Both counties show the same pattern: a huge
spike at 1 to 3 days (most patients go home quickly) and a long thin
tail going out to 30 days (a small fraction stay much longer). This
pattern is called "right-skewed" and is the reason this project uses
nonparametric tests instead of standard tests that assume a normal
bell curve. The two counties have nearly identical shapes, but Kings
has slightly more discharges overall.

\begin{figure}[H]
\centering
\includegraphics[width=0.95\textwidth]{figures/fig4_age_county.png}
\caption{Grouped bar plot of discharges by Age Group and County. The
70+ cohort is the single largest group in both counties.}
\label{fig:age-county}
\end{figure}

\textbf{What this plot shows.} This bar chart counts how many hospital
discharges fall into each age group, with Kings (blue) and Queens
(orange) shown side by side. Both counties follow the same pattern:
the elderly (70 or older) account for the most discharges, followed
by adults aged 50 to 69. The fewest discharges are for young adults
aged 18 to 29, who tend to be the healthiest age group. Kings County
has more discharges than Queens in every adult age group, which
matches its slightly larger total population.

\begin{figure}[H]
\centering
\includegraphics[width=0.7\textwidth]{figures/fig_pie_county.png}
\caption{Share of discharges between Kings (51.8\%, $n = 190{,}952$)
and Queens (48.2\%, $n = 177{,}543$) of the 368{,}495 total records.}
\label{fig:pie-county}
\end{figure}

\textbf{What this plot shows.} This pie chart shows what share of the
368,495 hospital discharges in this dataset came from each county.
Kings County (Brooklyn) accounts for slightly more than half (51.8\%)
and Queens accounts for the rest (48.2\%). The split is fairly even,
which means any group differences I find later in the paper come from
real differences in patient populations and not just from one county
having way more data than the other.

\begin{figure}[H]
\centering
\includegraphics[width=0.95\textwidth]{figures/fig5_ecdf_charges.png}
\caption{Empirical CDFs of Total Charges by County. The two CDFs
visibly separate, consistent with the Kolmogorov--Smirnov result
$D > 0$ and $p < 10^{-16}$.}
\label{fig:ecdf-charges}
\end{figure}

\textbf{What this plot shows.} An empirical CDF (cumulative
distribution function) shows what percentage of patients had charges
at or below a given dollar amount. Reading the chart at \$50,000 on
the x-axis, you can see that about 60\% of Queens patients are
charged at or below that amount, but only about 55\% of Kings
patients are. The Queens curve sits consistently above the Kings
curve, meaning Queens patients tend to have lower charges across the
board. The visible gap between the two curves is exactly what the
Kolmogorov-Smirnov test detects to confirm the distributions are
truly different.

\begin{figure}[H]
\centering
\includegraphics[width=0.85\textwidth]{figures/fig6_violin_severity.png}
\caption{Violin plot of $\log_{10}$(Total Costs) by All Patient
Refined Severity of Illness (1~=~Minor through 4~=~Extreme),
constructed from mirrored kernel densities with inner IQR boxes and
median dots. Both the center and spread of costs increase
monotonically from Minor to Extreme severity.}
\label{fig:violin-severity}
\end{figure}

\textbf{What this plot shows.} A violin plot is like a boxplot but
with the box "fattened" to show how common each cost level is. Each
violin represents one severity level, from Minor (light blue) on the
left to Extreme (dark blue) on the right. The wider sections show
where most patients in that severity group have their costs. As
severity worsens from Minor to Extreme, the violins shift upward (so
costs go up) and get wider at the top (so the most expensive cases
get even more expensive). This confirms that sicker patients cost the
hospital more, in a smooth and predictable way.

\begin{figure}[H]
\centering
\includegraphics[width=0.85\textwidth]{figures/fig7_age_charges.png}
\caption{Dot plot of median Total Charges by Age Group. There is a
strong monotone increase from the 0--17 cohort
(\textasciitilde\$20k) through the 70+ cohort
(\textasciitilde\$63k), consistent with the Kruskal--Wallis result
for Age Group.}
\label{fig:age-charges}
\end{figure}

\textbf{What this plot shows.} This dot plot shows the median (typical)
hospital charge for each age group, in thousands of dollars. The
pattern is very clear: the older the patient, the higher the typical
charge. A 70-year-old patient's typical hospital bill (about \$63,000)
is more than three times higher than a child's (about \$20,000). This
makes intuitive sense because older patients tend to have more complex
health conditions that need longer stays and more procedures. The
Kruskal-Wallis test confirms this trend is statistically significant.

% ============================================================
% 4. OBJECTIVES AND RESEARCH QUESTIONS (all 10 tests)
% ============================================================
\section{Objective Questions and Research Questions}

Table~\ref{tab:objectives} lists, separately for each (test, variable)
combination, the objective question (verbal form), the research
question (mathematical form), the null hypothesis, the test statistic,
the test name, the p-value, the decision at significance level
$\alpha=0.05$, and a short interpretation. This per-variable layout
makes it possible to read off, at a glance, exactly what was tested
on each individual outcome and exactly which records contributed to
each result. Throughout, $\alpha = 0.05$.

\begin{landscape}
\begin{scriptsize}
\renewcommand{\arraystretch}{1.15}
\begin{longtable}{p{2.0cm} p{2.4cm} p{3.0cm} p{2.7cm} p{2.5cm} p{1.7cm} p{1.3cm} p{1.3cm} p{3.0cm}}
\caption{Per-Variable Summary of All Ten Nonparametric Tests: Objective Question, Research Question, Null Hypothesis, Test Statistic, Test, P-value, Decision, and Interpretation}
\label{tab:objectives} \\
\toprule
\textbf{Test} & \textbf{Variable} & \textbf{Objective Question} & \textbf{Research Question} & \textbf{Null Hypothesis} & \textbf{Test Statistic} & \textbf{P-value} & \textbf{Decision} & \textbf{Interpretation} \\
\midrule
\endfirsthead
\multicolumn{9}{c}{\textit{(continued from previous page)}} \\
\toprule
\textbf{Test} & \textbf{Variable} & \textbf{Objective Question} & \textbf{Research Question} & \textbf{Null Hypothesis} & \textbf{Test Statistic} & \textbf{P-value} & \textbf{Decision} & \textbf{Interpretation} \\
\midrule
\endhead

%====================== RUNS TEST (7 rows) ======================
\multicolumn{9}{l}{\textit{\textbf{Runs Test (7 testable variables)}}} \\

Runs Test & Length of Stay & Is Length of Stay recorded in random order across the dataset? & Does the observed number of runs in the LOS sequence match the expected number under randomness? & $H_0$: the LOS sequence is random. & $Z = -56.17$ & $< 2.2 \times 10^{-16}$ & Reject $H_0$ & LOS values cluster by facility and admission date rather than appearing randomly. \\

Runs Test & Total Charges & Are Total Charges recorded in random order across the dataset? & Does the observed number of runs in the Charges sequence match the expected number under randomness? & $H_0$: the Charges sequence is random. & $Z = -90.68$ & $< 2.2 \times 10^{-16}$ & Reject $H_0$ & Charges are clustered by facility, so adjacent records have similar billing patterns. \\

Runs Test & Total Costs & Are Total Costs recorded in random order across the dataset? & Does the observed number of runs in the Costs sequence match the expected number under randomness? & $H_0$: the Costs sequence is random. & $Z = -75.34$ & $< 2.2 \times 10^{-16}$ & Reject $H_0$ & Costs are clustered by facility, mirroring the pattern seen in Charges. \\

Runs Test & Birth Weight & Are Birth Weight records (newborns only) in random order? & Does the observed number of runs in the Birth Weight sequence match the expected number under randomness? & $H_0$: the Birth Weight sequence is random. & $Z = -1.98$ & $0.0477$ & Reject $H_0$ & Borderline rejection; ordering is only weakly non-random. \\

Runs Test & APR DRG Code & Is the All Patient Refined Diagnosis-Related Group label in random order? & Does the observed number of runs match the expected number under randomness? & $H_0$: the DRG-code sequence is random. & $Z = -53.29$ & $< 2.2 \times 10^{-16}$ & Reject $H_0$ & Records sorted by facility tend to share DRG codes; reflects file ordering, not a clinical pattern. \\

Runs Test & APR MDC Code & Is the All Patient Refined Major Diagnostic Category label in random order? & Does the observed number of runs match the expected number under randomness? & $H_0$: the MDC-code sequence is random. & $Z = -54.02$ & $< 2.2 \times 10^{-16}$ & Reject $H_0$ & Same explanation as for DRG: facilities specialise in subsets of MDCs, so the codes cluster. \\

Runs Test & APR Severity & Is the APR Severity of Illness Code in random order? & Does the observed number of runs match the expected number under randomness? & $H_0$: the Severity-code sequence is random. & $Z = -33.12$ & $< 2.2 \times 10^{-16}$ & Reject $H_0$ & Severity codes cluster by facility because hospitals see different case-mixes. \\

%====================== SIGN TEST (5 rows) ======================
\midrule
\multicolumn{9}{l}{\textit{\textbf{Sign Test (5 paired variables: Kings vs Queens)}}} \\

Sign Test & Length of Stay & Is the median Length of Stay the same in Kings and Queens patients? & Does the median of paired LOS differences (Kings $-$ Queens) equal zero? & $H_0$: $\widetilde{\mathrm{LOS}}_K = \widetilde{\mathrm{LOS}}_Q$. & $S = 81{,}301$ & $< 2.2 \times 10^{-16}$ & Reject $H_0$ & Kings patients tend to have longer stays than Queens patients. \\

Sign Test & Total Charges & Is the median Total Charge the same in Kings and Queens patients? & Does the median of paired Charges differences (Kings $-$ Queens) equal zero? & $H_0$: $\widetilde{\mathrm{Charges}}_K = \widetilde{\mathrm{Charges}}_Q$. & $S = 94{,}864$ & $< 2.2 \times 10^{-16}$ & Reject $H_0$ & Median Charges in Kings exceed those in Queens. \\

Sign Test & Total Costs & Is the median Total Cost the same in Kings and Queens admissions? & Does the median of paired Costs differences (Kings $-$ Queens) equal zero? & $H_0$: $\widetilde{\mathrm{Costs}}_K = \widetilde{\mathrm{Costs}}_Q$. & $S = 105{,}522$ & $< 2.2 \times 10^{-16}$ & Reject $H_0$ & Median Costs in Kings exceed those in Queens, mirroring the Charges result. \\

Sign Test & Birth Weight & Is the median Birth Weight the same in Kings and Queens newborns? & Does the median of paired Birth Weight differences (Kings $-$ Queens) equal zero? & $H_0$: $\widetilde{\mathrm{BW}}_K = \widetilde{\mathrm{BW}}_Q$. & $S = 10{,}374$ & $< 2.2 \times 10^{-16}$ & Reject $H_0$ & Median Birth Weight differs between newborns delivered in Kings vs Queens. \\

Sign Test & APR Severity & Is the median APR Severity of Illness Code the same in Kings and Queens patients? & Does the median of paired Severity differences (Kings $-$ Queens) equal zero? & $H_0$: $\widetilde{\mathrm{Sev}}_K = \widetilde{\mathrm{Sev}}_Q$. & $S = 65{,}885$ & $< 2.2 \times 10^{-16}$ & Reject $H_0$ & Severity case-mix differs systematically between the two counties. \\

%====================== WILCOXON SIGNED-RANK (5 rows) ======================
\midrule
\multicolumn{9}{l}{\textit{\textbf{Wilcoxon Signed-Rank (5 paired variables: Kings vs Queens)}}} \\

Wilcoxon Signed-Rank & Length of Stay & Do paired LOS distributions differ between counties, accounting for the size of differences? & Is the median of signed-rank LOS differences zero? & $H_0$: median of signed-rank LOS differences $= 0$. & $V = 6.40 \times 10^{9}$ & $< 2.2 \times 10^{-16}$ & Reject $H_0$ & LOS differs between counties even when difference magnitudes are accounted for. \\

Wilcoxon Signed-Rank & Total Charges & Do paired Charges distributions differ between counties, accounting for the size of differences? & Is the median of signed-rank Charges differences zero? & $H_0$: median of signed-rank Charges differences $= 0$. & $V = 8.59 \times 10^{9}$ & $< 2.2 \times 10^{-16}$ & Reject $H_0$ & Charges differ between counties; the magnitude of the difference is also non-trivial. \\

Wilcoxon Signed-Rank & Total Costs & Do paired Costs distributions differ between counties, accounting for the size of differences? & Is the median of signed-rank Costs differences zero? & $H_0$: median of signed-rank Costs differences $= 0$. & $V = 9.92 \times 10^{9}$ & $< 2.2 \times 10^{-16}$ & Reject $H_0$ & Costs distributions differ between counties beyond a simple location shift. \\

Wilcoxon Signed-Rank & Birth Weight & Do paired Birth Weight distributions differ between counties, accounting for the size of differences? & Is the median of signed-rank Birth Weight differences zero? & $H_0$: median of signed-rank Birth Weight differences $= 0$. & $V = 1.04 \times 10^{8}$ & $< 2.2 \times 10^{-16}$ & Reject $H_0$ & Newborn Birth Weight distributions differ between Kings and Queens. \\

Wilcoxon Signed-Rank & APR Severity & Do paired Severity distributions differ between counties, accounting for the size of differences? & Is the median of signed-rank Severity differences zero? & $H_0$: median of signed-rank Severity differences $= 0$. & $V = 4.17 \times 10^{9}$ & $< 2.2 \times 10^{-16}$ & Reject $H_0$ & Severity distributions differ between counties beyond a simple median shift. \\

%====================== MANN-WHITNEY U (6 rows) ======================
\midrule
\multicolumn{9}{l}{\textit{\textbf{Mann-Whitney U / Wilcoxon Rank Sum (6 independent-sample comparisons)}}} \\

Mann-Whitney U & Length of Stay (Kings vs Queens) & Do independent-sample LOS distributions differ between Kings and Queens? & Are the rank-sums of LOS in the two groups equal? & $H_0$: $F_{\mathrm{LOS}}^{K} = F_{\mathrm{LOS}}^{Q}$. & $W = 1.76 \times 10^{10}$ & $< 2.2 \times 10^{-16}$ & Reject $H_0$ & LOS distributions differ between counties as independent samples. \\

Mann-Whitney U & Total Charges (Kings vs Queens) & Do independent-sample Charges distributions differ between Kings and Queens? & Are the rank-sums of Charges in the two groups equal? & $H_0$: $F_{\mathrm{Charges}}^{K} = F_{\mathrm{Charges}}^{Q}$. & $W = 1.81 \times 10^{10}$ & $< 2.2 \times 10^{-16}$ & Reject $H_0$ & Charges distributions differ between counties. \\

Mann-Whitney U & Total Costs (Kings vs Queens) & Do independent-sample Costs distributions differ between Kings and Queens? & Are the rank-sums of Costs in the two groups equal? & $H_0$: $F_{\mathrm{Costs}}^{K} = F_{\mathrm{Costs}}^{Q}$. & $W = 2.02 \times 10^{10}$ & $< 2.2 \times 10^{-16}$ & Reject $H_0$ & Costs distributions differ between counties. \\

Mann-Whitney U & Birth Weight (Kings vs Queens) & Do independent-sample Birth Weight distributions differ between Kings and Queens? & Are the rank-sums of Birth Weight equal in the two groups? & $H_0$: $F_{\mathrm{BW}}^{K} = F_{\mathrm{BW}}^{Q}$. & $W = 2.49 \times 10^{8}$ & $< 2.2 \times 10^{-16}$ & Reject $H_0$ & Birth Weight distributions differ between Kings and Queens newborns. \\

Mann-Whitney U & APR Severity (Kings vs Queens) & Do independent-sample Severity distributions differ between Kings and Queens? & Are the rank-sums of Severity equal in the two groups? & $H_0$: $F_{\mathrm{Sev}}^{K} = F_{\mathrm{Sev}}^{Q}$. & $W = 1.75 \times 10^{10}$ & $< 2.2 \times 10^{-16}$ & Reject $H_0$ & Severity case-mix differs between counties as independent samples. \\

Mann-Whitney U & Total Charges (ED vs Non-ED) & Do Charges differ between Emergency Department and Non-Emergency Department admissions? & Are rank-sums of Charges equal across the two admission modes? & $H_0$: $F_{\mathrm{Charges}}^{\mathrm{ED}} = F_{\mathrm{Charges}}^{\mathrm{Non\text{-}ED}}$. & $W = 1.08 \times 10^{10}$ & $< 2.2 \times 10^{-16}$ & Reject $H_0$ & ED admissions have a different cost profile than non-ED admissions. \\

%====================== KOLMOGOROV-SMIRNOV (8 rows) ======================
\midrule
\multicolumn{9}{l}{\textit{\textbf{Kolmogorov-Smirnov Two-Sample (4 outcomes $\times$ 2 grouping schemes = 8 comparisons)}}} \\

Kolmogorov-Smirnov & Length of Stay (Kings vs Queens) & Does the entire CDF of LOS differ between Kings and Queens? & Is $F_{K}(x) = F_{Q}(x)$ for all $x$? & $H_0$: CDFs of LOS are identical in the two counties. & $D = 0.0298$ & $< 2.2 \times 10^{-16}$ & Reject $H_0$ & LOS distribution shapes, not just centers, differ between counties. \\

Kolmogorov-Smirnov & Total Charges (Kings vs Queens) & Does the entire CDF of Charges differ between Kings and Queens? & Is $F_{K}(x) = F_{Q}(x)$ for all $x$? & $H_0$: CDFs of Charges are identical in the two counties. & $D = 0.0551$ & $< 2.2 \times 10^{-16}$ & Reject $H_0$ & Charges distribution shapes differ across the entire dollar range. \\

Kolmogorov-Smirnov & Total Costs (Kings vs Queens) & Does the entire CDF of Costs differ between Kings and Queens? & Is $F_{K}(x) = F_{Q}(x)$ for all $x$? & $H_0$: CDFs of Costs are identical in the two counties. & $D = 0.1439$ & $< 2.2 \times 10^{-16}$ & Reject $H_0$ & Costs show the largest county-level shape difference among the four outcomes. \\

Kolmogorov-Smirnov & Birth Weight (Kings vs Queens) & Does the entire CDF of Birth Weight differ between Kings and Queens? & Is $F_{K}(x) = F_{Q}(x)$ for all $x$? & $H_0$: CDFs of Birth Weight are identical in the two counties. & $D = 0.0464$ & $< 2.2 \times 10^{-16}$ & Reject $H_0$ & Birth Weight distribution shapes differ between Kings and Queens newborns. \\

Kolmogorov-Smirnov & Length of Stay (ED vs Non-ED) & Does the entire CDF of LOS differ between Emergency and Non-Emergency admissions? & Is $F_{\mathrm{ED}}(x) = F_{\overline{\mathrm{ED}}}(x)$ for all $x$? & $H_0$: CDFs of LOS are identical for ED and non-ED admissions. & $D = 0.2959$ & $< 2.2 \times 10^{-16}$ & Reject $H_0$ & Largest CDF separation in the table; ED stays differ profoundly from non-ED stays. \\

Kolmogorov-Smirnov & Total Charges (ED vs Non-ED) & Does the entire CDF of Charges differ between Emergency and Non-Emergency admissions? & Is $F_{\mathrm{ED}}(x) = F_{\overline{\mathrm{ED}}}(x)$ for all $x$? & $H_0$: CDFs of Charges are identical for ED and non-ED admissions. & $D = 0.1950$ & $< 2.2 \times 10^{-16}$ & Reject $H_0$ & ED billing patterns differ from non-ED billing patterns. \\

Kolmogorov-Smirnov & Total Costs (ED vs Non-ED) & Does the entire CDF of Costs differ between Emergency and Non-Emergency admissions? & Is $F_{\mathrm{ED}}(x) = F_{\overline{\mathrm{ED}}}(x)$ for all $x$? & $H_0$: CDFs of Costs are identical for ED and non-ED admissions. & $D = 0.1674$ & $< 2.2 \times 10^{-16}$ & Reject $H_0$ & ED admissions cost the hospital differently than scheduled admissions. \\

Kolmogorov-Smirnov & Birth Weight (ED vs Non-ED) & Does the entire CDF of Birth Weight differ between Emergency and Non-Emergency admissions? & Is $F_{\mathrm{ED}}(x) = F_{\overline{\mathrm{ED}}}(x)$ for all $x$? & $H_0$: CDFs of Birth Weight are identical for ED and non-ED admissions. & $D = 0.0530$ & $0.00155$ & Reject $H_0$ & Smallest of the eight effects but still significant. \\

%====================== KRUSKAL-WALLIS (15 rows) ======================
\midrule
\multicolumn{9}{l}{\textit{\textbf{Kruskal-Wallis (4 outcomes $\times$ 4 groupings = 15 testable combinations; BirthWt $\times$ Age Group is degenerate)}}} \\

Kruskal-Wallis & Length of Stay $\times$ Age Group & Does the LOS distribution differ across age groups? & Are LOS distributions identical across all five age groups? & $H_0$: $F_1 = F_2 = \cdots = F_5$. & $\chi^2 = 39{,}479$ ($df=4$) & $< 2.2 \times 10^{-16}$ & Reject $H_0$ & LOS varies strongly with patient age. \\

Kruskal-Wallis & Length of Stay $\times$ APR Severity & Does the LOS distribution differ across severity levels? & Are LOS distributions identical across the four severity levels? & $H_0$: $F_1 = F_2 = F_3 = F_4$. & $\chi^2 = 82{,}402$ ($df=4$) & $< 2.2 \times 10^{-16}$ & Reject $H_0$ & Severity is the strongest single driver of LOS in this dataset. \\

Kruskal-Wallis & Length of Stay $\times$ Type of Admission & Does the LOS distribution differ across admission types? & Are LOS distributions identical across admission types? & $H_0$: $F_1 = F_2 = \cdots = F_6$. & $\chi^2 = 25{,}198$ ($df=5$) & $< 2.2 \times 10^{-16}$ & Reject $H_0$ & Elective, emergency, urgent, newborn, and trauma admissions have different LOS profiles. \\

Kruskal-Wallis & Length of Stay $\times$ Primary Payer & Does the LOS distribution differ across primary payer categories? & Are LOS distributions identical across payer categories? & $H_0$: all payer-group LOS distributions are equal. & $\chi^2 = 28{,}581$ ($df=8$) & $< 2.2 \times 10^{-16}$ & Reject $H_0$ & LOS depends on payer (Medicare, Medicaid, Private, etc.). \\

Kruskal-Wallis & Total Charges $\times$ Age Group & Does the Charges distribution differ across age groups? & Are Charges distributions identical across age groups? & $H_0$: all age-group Charges distributions are equal. & $\chi^2 = 65{,}223$ ($df=4$) & $< 2.2 \times 10^{-16}$ & Reject $H_0$ & Charges rise sharply with age (see Figure~\ref{fig:age-charges}). \\

Kruskal-Wallis & Total Charges $\times$ APR Severity & Does the Charges distribution differ across severity levels? & Are Charges distributions identical across the four severity levels? & $H_0$: all severity-group Charges distributions are equal. & $\chi^2 = 78{,}528$ ($df=4$) & $< 2.2 \times 10^{-16}$ & Reject $H_0$ & Sicker patients are billed substantially more. \\

Kruskal-Wallis & Total Charges $\times$ Type of Admission & Does the Charges distribution differ across admission types? & Are Charges distributions identical across admission types? & $H_0$: all admission-type Charges distributions are equal. & $\chi^2 = 45{,}558$ ($df=5$) & $< 2.2 \times 10^{-16}$ & Reject $H_0$ & Charges depend on whether admission was elective, urgent, etc. \\

Kruskal-Wallis & Total Charges $\times$ Primary Payer & Does the Charges distribution differ across primary payer categories? & Are Charges distributions identical across payer categories? & $H_0$: all payer-group Charges distributions are equal. & $\chi^2 = 35{,}087$ ($df=8$) & $< 2.2 \times 10^{-16}$ & Reject $H_0$ & Charges vary by payer mix. \\

Kruskal-Wallis & Total Costs $\times$ Age Group & Does the Costs distribution differ across age groups? & Are Costs distributions identical across age groups? & $H_0$: all age-group Costs distributions are equal. & $\chi^2 = 61{,}785$ ($df=4$) & $< 2.2 \times 10^{-16}$ & Reject $H_0$ & Hospital Costs rise with age, mirroring the Charges pattern. \\

Kruskal-Wallis & Total Costs $\times$ APR Severity & Does the Costs distribution differ across severity levels? & Are Costs distributions identical across severity levels? & $H_0$: all severity-group Costs distributions are equal. & $\chi^2 = 76{,}947$ ($df=4$) & $< 2.2 \times 10^{-16}$ & Reject $H_0$ & Severity drives Costs as strongly as it drives Charges. \\

Kruskal-Wallis & Total Costs $\times$ Type of Admission & Does the Costs distribution differ across admission types? & Are Costs distributions identical across admission types? & $H_0$: all admission-type Costs distributions are equal. & $\chi^2 = 50{,}005$ ($df=5$) & $< 2.2 \times 10^{-16}$ & Reject $H_0$ & Costs depend on the admission pathway. \\

Kruskal-Wallis & Total Costs $\times$ Primary Payer & Does the Costs distribution differ across primary payer categories? & Are Costs distributions identical across payer categories? & $H_0$: all payer-group Costs distributions are equal. & $\chi^2 = 25{,}578$ ($df=8$) & $< 2.2 \times 10^{-16}$ & Reject $H_0$ & Costs vary by payer mix. \\

Kruskal-Wallis & Birth Weight $\times$ APR Severity & Does newborn Birth Weight differ across severity levels? & Are Birth Weight distributions identical across severity levels? & $H_0$: all severity-group Birth Weight distributions are equal. & $\chi^2 = 1{,}561.7$ ($df=4$) & $< 2.2 \times 10^{-16}$ & Reject $H_0$ & Higher-severity newborns tend to have lower Birth Weights. \\

Kruskal-Wallis & Birth Weight $\times$ Type of Admission & Does newborn Birth Weight differ across admission types? & Are Birth Weight distributions identical across admission types? & $H_0$: all admission-type Birth Weight distributions are equal. & $\chi^2 = 40.5$ ($df=4$) & $< 2.2 \times 10^{-16}$ & Reject $H_0$ & Birth Weight differs across admission pathways, but more weakly. \\

Kruskal-Wallis & Birth Weight $\times$ Primary Payer & Does newborn Birth Weight differ across primary payer categories? & Are Birth Weight distributions identical across payer categories? & $H_0$: all payer-group Birth Weight distributions are equal. & $\chi^2 = 52.6$ ($df=7$) & $< 2.2 \times 10^{-16}$ & Reject $H_0$ & Birth Weight is mildly associated with payer category. \\

%====================== FRIEDMAN (3 rows) ======================
\midrule
\multicolumn{9}{l}{\textit{\textbf{Friedman Test (Severity as treatment, Facility as block; 3 outcomes)}}} \\

Friedman & Length of Stay & Within each facility, does illness severity change LOS? & After blocking by facility, is there a severity effect on LOS? & $H_0$: no severity (treatment) effect on LOS within blocks. & $Q = 13.29$ ($df=3$) & $0.0041$ & Reject $H_0$ & Even after removing facility-level differences, severity drives LOS. \\

Friedman & Total Charges & Within each facility, does illness severity change Charges? & After blocking by facility, is there a severity effect on Charges? & $H_0$: no severity (treatment) effect on Charges within blocks. & $Q = 15.00$ ($df=3$) & $0.0018$ & Reject $H_0$ & Severity drives Charges within facility, isolating clinical from pricing variation. \\

Friedman & Total Costs & Within each facility, does illness severity change Costs? & After blocking by facility, is there a severity effect on Costs? & $H_0$: no severity (treatment) effect on Costs within blocks. & $Q = 15.00$ ($df=3$) & $0.0018$ & Reject $H_0$ & Severity drives Costs within facility; the same finding as for Charges. \\

%====================== KENDALL (6 rows) ======================
\midrule
\multicolumn{9}{l}{\textit{\textbf{Kendall's $\tau$ (6 numeric pairs; subsample $n=5{,}000$)}}} \\

Kendall's $\tau$ & LOS vs Charges & Is there an ordinal association between LOS and Charges? & Is Kendall's $\tau$ between LOS and Charges zero? & $H_0$: $\tau = 0$. & $\tau = 0.5761$ & $< 2.2 \times 10^{-16}$ & Reject $H_0$ & LOS and Charges are strongly concordant: longer stays cost more. \\

Kendall's $\tau$ & LOS vs Costs & Is there an ordinal association between LOS and Costs? & Is Kendall's $\tau$ between LOS and Costs zero? & $H_0$: $\tau = 0$. & $\tau = 0.5858$ & $< 2.2 \times 10^{-16}$ & Reject $H_0$ & Longer stays raise actual hospital costs, in close lockstep with Charges. \\

Kendall's $\tau$ & LOS vs Birth Weight & Is there an ordinal association between LOS and Birth Weight? & Is Kendall's $\tau$ between LOS and Birth Weight zero? & $H_0$: $\tau = 0$. & $\tau = -0.1618$ & $7.0 \times 10^{-7}$ & Reject $H_0$ & Lower-Birth-Weight newborns tend to have longer stays (weak negative). \\

Kendall's $\tau$ & Charges vs Costs & Is there an ordinal association between Charges and Costs? & Is Kendall's $\tau$ between Charges and Costs zero? & $H_0$: $\tau = 0$. & $\tau = 0.6708$ & $< 2.2 \times 10^{-16}$ & Reject $H_0$ & Charges and Costs are strongly concordant; consistent markup ratios. \\

Kendall's $\tau$ & Charges vs Birth Weight & Is there an ordinal association between Charges and Birth Weight? & Is Kendall's $\tau$ between Charges and Birth Weight zero? & $H_0$: $\tau = 0$. & $\tau = -0.1145$ & $6.9 \times 10^{-5}$ & Reject $H_0$ & Lower Birth Weight is mildly associated with higher Charges. \\

Kendall's $\tau$ & Costs vs Birth Weight & Is there an ordinal association between Costs and Birth Weight? & Is Kendall's $\tau$ between Costs and Birth Weight zero? & $H_0$: $\tau = 0$. & $\tau = -0.1100$ & $1.3 \times 10^{-4}$ & Reject $H_0$ & Lower Birth Weight is mildly associated with higher Costs. \\

%====================== SPEARMAN (6 rows) ======================
\midrule
\multicolumn{9}{l}{\textit{\textbf{Spearman's $\rho$ (6 numeric pairs; full data)}}} \\

Spearman's $\rho$ & LOS vs Charges & Is there a monotone rank association between LOS and Charges? & Is Spearman's $\rho$ between LOS and Charges zero? & $H_0$: $\rho_s = 0$. & $\rho_s = 0.7305$ & $< 2.2 \times 10^{-16}$ & Reject $H_0$ & LOS and Charges are tightly coupled in rank. \\

Spearman's $\rho$ & LOS vs Costs & Is there a monotone rank association between LOS and Costs? & Is Spearman's $\rho$ between LOS and Costs zero? & $H_0$: $\rho_s = 0$. & $\rho_s = 0.7478$ & $< 2.2 \times 10^{-16}$ & Reject $H_0$ & LOS tracks Costs even more tightly than it tracks Charges. \\

Spearman's $\rho$ & LOS vs Birth Weight & Is there a monotone rank association between LOS and Birth Weight? & Is Spearman's $\rho$ between LOS and Birth Weight zero? & $H_0$: $\rho_s = 0$. & $\rho_s = -0.1733$ & $< 2.2 \times 10^{-16}$ & Reject $H_0$ & Heavier newborns have shorter stays; weak but significant. \\

Spearman's $\rho$ & Charges vs Costs & Is there a monotone rank association between Charges and Costs? & Is Spearman's $\rho$ between Charges and Costs zero? & $H_0$: $\rho_s = 0$. & $\rho_s = 0.8589$ & $< 2.2 \times 10^{-16}$ & Reject $H_0$ & Strongest pair in the table; Charges and Costs are nearly proportional in rank. \\

Spearman's $\rho$ & Charges vs Birth Weight & Is there a monotone rank association between Charges and Birth Weight? & Is Spearman's $\rho$ between Charges and Birth Weight zero? & $H_0$: $\rho_s = 0$. & $\rho_s = -0.1257$ & $< 2.2 \times 10^{-16}$ & Reject $H_0$ & Heavier newborns are billed less; weak but significant. \\

Spearman's $\rho$ & Costs vs Birth Weight & Is there a monotone rank association between Costs and Birth Weight? & Is Spearman's $\rho$ between Costs and Birth Weight zero? & $H_0$: $\rho_s = 0$. & $\rho_s = -0.1490$ & $< 2.2 \times 10^{-16}$ & Reject $H_0$ & Heavier newborns cost the hospital less; weak but significant. \\

%====================== PERMUTATION (1 row) ======================
\midrule
\multicolumn{9}{l}{\textit{\textbf{Permutation Test (Monte Carlo, B = 5{,}000)}}} \\

Permutation & Total Charges (Kings vs Queens) & After randomly shuffling county labels, is the observed median Charges difference more extreme than chance? & Is the median Charges difference between counties attributable to chance under exchangeability? & $H_0$: county label is exchangeable; the median of Charges is the same in both counties. & Observed median diff $= \$2{,}858.07$ & $0.011$ & Reject $H_0$ & The median Charges in Kings and Queens are significantly different. The robust median detects a real distributional shift even with extreme outliers in Kings. \\

\bottomrule
\end{longtable}
\end{scriptsize}
\end{landscape}

% ============================================================
% 5. NONPARAMETRIC ANALYSES (all 10 tests)
% ============================================================
\section{Nonparametric Analyses}

\subsection{Test 1: Runs Test for Randomness}
\label{sec:runs}

The Runs Test is a nonparametric test used to assess whether a sequence
of observations follows a random pattern. Each quantitative variable
was classified above or below its median, and the number of runs
(consecutive sequences of the same category) was counted.

\paragraph{Test statistic.} For a sequence dichotomised around its
median, with $n_1$ ones and $n_2$ zeros ($n=n_1+n_2$), let $R$ be the
observed number of runs. Under $H_0$,
\[
  E(R) = \frac{2n_1 n_2}{n} + 1,\qquad
  \operatorname{Var}(R) = \frac{2n_1 n_2 (2n_1 n_2 - n)}{n^2 (n-1)},\qquad
  Z = \frac{R - E(R)}{\sqrt{\operatorname{Var}(R)}}.
\]

\begin{table}[H]
\centering
\caption{Runs Test Results, Statewide Planning and Research Cooperative System 2022 K\&Q}
\label{tab:runs}
\begin{tabular}{lccr l}
\toprule
\textbf{Variable} & $R$ (obs.) & $E(R)$ & \textbf{$p$-value} & \textbf{Decision} \\
\midrule
Length of Stay                & 164{,}962 & 181{,}791 & $< 2.2\times 10^{-16}$ & Reject $H_0$ \\
Total Charges                 & 156{,}725 & 184{,}249 & $< 2.2\times 10^{-16}$ & Reject $H_0$ \\
Total Costs                   & 161{,}381 & 184{,}249 & $< 2.2\times 10^{-16}$ & Reject $H_0$ \\
Birth Weight                  &  21{,}248 &  21{,}452 & $0.0477$              & Reject $H_0$ \\
All Patient Refined Diagnosis-Related Group Code                  & 168{,}061 & 184{,}233 & $< 2.2\times 10^{-16}$ & Reject $H_0$ \\
All Patient Refined Major Diagnostic Category Code                  & 167{,}813 & 184{,}206 & $< 2.2\times 10^{-16}$ & Reject $H_0$ \\
All Patient Refined Severity of Illness Code  & 152{,}218 & 161{,}003 & $< 2.2\times 10^{-16}$ & Reject $H_0$ \\
Discharge Year                & \multicolumn{3}{c}{Not testable (constant = 2022)} & --- \\
\bottomrule
\end{tabular}
\end{table}

\paragraph{Interpretation.}
At $\alpha = 0.05$, every testable quantitative variable turned out to
be non-random. Honestly, this is not too surprising. Hospital discharge
data is organized by facility and by when patients are admitted, so
there is going to be some natural structure in the ordering. Birth
Weight was borderline ($p = 0.0477$), just barely under the cutoff.
Discharge Year could not be tested at all since every record is from
2022, so there is no variability to work with. This result is
consistent with Bujang and Sapri (2018), who found that
consecutive-sampling clinical data almost always shows non-random
ordering.

\paragraph{Limitation note on nominal codes.}
The Runs Test technically requires an ordered numeric sequence so that
``above the median'' and ``below the median'' carry meaning. Two of the
seven variables tested above --- the All Patient Refined Diagnosis-Related
Group Code and the All Patient Refined Major Diagnostic Category Code ---
are nominal classification labels stored as integers, so dichotomising
them around their numeric median is mathematically valid but not
clinically interpretable. The rejection of $H_0$ for those two rows
reflects only that the file is sorted by facility (each facility tends
to specialise in a subset of Diagnosis-Related Group and Major
Diagnostic Category codes), not a clinical pattern. The four genuinely
quantitative variables (Length of Stay, Total Charges, Total Costs,
Birth Weight) plus the ordinal All Patient Refined Severity of Illness
Code give the substantively meaningful Runs Test results.

\subsection{Test 2: Sign Test}

The Sign Test is a nonparametric test for paired samples that examines
whether the median difference between two related groups is zero. It
only considers the direction (sign) of each difference.

\paragraph{Methodology.} For paired Length of Stay observations between
Kings ($K_i$) and Queens ($Q_i$):
\begin{enumerate}[nosep]
    \item Compute differences $D_i = K_i - Q_i$.
    \item Remove zero differences.
    \item Count positive differences $n^+$ and negative differences $n^-$.
    \item Under $H_0$, $n^+ \sim \text{Binomial}(n, 0.5)$.
\end{enumerate}

\paragraph{Results.}
The dependent-samples Sign Test was applied to five clinically
meaningful quantitative variables. Detailed results for Length of Stay
(the flagship paired comparison):
\[
S = 81{,}301, \quad p < 2.2 \times 10^{-16}.
\]

Table~\ref{tab:sign-all} presents the Sign Test results across all
five quantitative variables for the Kings-vs-Queens paired comparison.

\begin{table}[H]
\centering
\caption{Sign Test Results: Kings vs Queens (Paired)}
\label{tab:sign-all}
\begin{tabular}{lcc}
\toprule
\textbf{Variable ($H_0$: $\tilde{\mu}_K = \tilde{\mu}_Q$)} & \textbf{$p$-value} & \textbf{Decision} \\
\midrule
Length of Stay (days)           & $< 2.2 \times 10^{-16}$ & Reject $H_0$ \\
Total Charges (US Dollars)             & $< 2.2 \times 10^{-16}$ & Reject $H_0$ \\
Total Costs (US Dollars)               & $< 2.2 \times 10^{-16}$ & Reject $H_0$ \\
Birth Weight (grams)            & $< 2.2 \times 10^{-16}$ & Reject $H_0$ \\
All Patient Refined Severity of Illness Code    & $< 2.2 \times 10^{-16}$ & Reject $H_0$ \\
\bottomrule
\end{tabular}
\end{table}

\paragraph{Interpretation.}
With p-values way below 0.05 for every comparison, we reject the null
hypothesis in all five cases. The data gives us strong evidence that
the median Length of Stay is different between Kings and Queens County
hospitals---from the direction of the differences, Kings County
patients tend to have longer stays than those in Queens. The same
conclusion holds for Total Charges, Total Costs, Birth Weight (for
newborn discharges), and All Patient Refined Severity of Illness Code: every one of
these medians differs significantly between the two counties.

\subsection{Test 3: Wilcoxon Signed-Rank Test}

The Wilcoxon Signed-Rank Test extends the Sign Test by incorporating the
magnitude of differences, not just their direction.

\paragraph{Methodology.}
\begin{enumerate}[nosep]
    \item Compute paired differences $D_i$.
    \item Rank the absolute differences $|D_i|$.
    \item Sum the ranks of positive differences ($T^+$) and negative differences ($T^-$).
    \item The test statistic is $V = \min(T^+, T^-)$.
\end{enumerate}

\paragraph{Results.}

Detailed result for Length of Stay (Kings vs Queens):
\[
V = 6{,}399{,}613{,}688, \quad p < 2.2 \times 10^{-16}
\]

Table~\ref{tab:wsr-all} presents the Wilcoxon Signed-Rank results
across all five quantitative variables for the Kings-vs-Queens paired
comparison.

\begin{table}[H]
\centering
\caption{Wilcoxon Signed-Rank Test Results: Kings vs Queens (Paired)}
\label{tab:wsr-all}
\begin{tabular}{lcc}
\toprule
\textbf{Variable ($H_0$: median paired difference $= 0$)} & \textbf{$p$-value} & \textbf{Decision} \\
\midrule
Length of Stay (days)           & $< 2.2 \times 10^{-16}$ & Reject $H_0$ \\
Total Charges (US Dollars)             & $< 2.2 \times 10^{-16}$ & Reject $H_0$ \\
Total Costs (US Dollars)               & $< 2.2 \times 10^{-16}$ & Reject $H_0$ \\
Birth Weight (grams)            & $< 2.2 \times 10^{-16}$ & Reject $H_0$ \\
All Patient Refined Severity of Illness Code    & $< 2.2 \times 10^{-16}$ & Reject $H_0$ \\
\bottomrule
\end{tabular}
\end{table}

\paragraph{Interpretation.}
What is nice about the Wilcoxon Signed-Rank Test compared to the Sign
Test is that it does not just look at which direction the difference
goes---it also takes into account how large the difference is. The
results here back up what the Sign Test found, but with even stronger
evidence. All five variables---Length of Stay, Total Charges, Total
Costs, Birth Weight, and All Patient Refined Severity of Illness---show significant
distributional differences between Kings and Queens counties.

\paragraph{Combined paired tests summary.}
Table~\ref{tab:paired-summary} consolidates the Sign Test and Wilcoxon
Signed-Rank Test results side-by-side for easy comparison, covering
all five quantitative variables tested as paired Kings-vs-Queens
comparisons. ID numbers and classification codes (Operating
Certificate Number, Permanent Facility ID, All Patient Refined Diagnosis-Related Group Code, All Patient Refined Major Diagnostic Category Code)
are excluded because comparing their medians between counties has no
practical meaning. Discharge Year is constant (2022) and cannot be
tested.

\begin{table}[H]
\centering
\caption{Combined Paired Tests Summary: Sign Test and Wilcoxon Signed-Rank}
\label{tab:paired-summary}
\small
\begin{tabular}{p{0.30\textwidth} c c c c}
\toprule
 & \multicolumn{2}{c}{\textbf{P-value}} & \multicolumn{2}{c}{\textbf{Decision ($\alpha = 0.05$)}} \\
\cmidrule(lr){2-3} \cmidrule(lr){4-5}
\textbf{$H_0$: $\tilde{\mu}_K = \tilde{\mu}_Q$} & \textbf{Sign Test} & \textbf{W. Signed-Rank} & \textbf{Sign} & \textbf{Wilcoxon Signed-Rank} \\
\midrule
Length of Stay (days)        & $< 2.2\times 10^{-16}$ & $< 2.2\times 10^{-16}$ & Reject & Reject \\
Total Charges (US Dollars)          & $< 2.2\times 10^{-16}$ & $< 2.2\times 10^{-16}$ & Reject & Reject \\
Total Costs (US Dollars)            & $< 2.2\times 10^{-16}$ & $< 2.2\times 10^{-16}$ & Reject & Reject \\
Birth Weight (grams)         & $< 2.2\times 10^{-16}$ & $< 2.2\times 10^{-16}$ & Reject & Reject \\
All Patient Refined Severity of Illness      & $< 2.2\times 10^{-16}$ & $< 2.2\times 10^{-16}$ & Reject & Reject \\
\bottomrule
\end{tabular}
\end{table}

\subsection{Test 4: Wilcoxon Rank Sum Test (Mann--Whitney U Test)}

The Wilcoxon Rank Sum Test (equivalent to the Mann--Whitney U Test)
compares two independent groups without requiring the assumption of
normality.

\paragraph{Methodology.}
\begin{enumerate}[nosep]
    \item Combine all observations and rank them from smallest to largest.
    \item Compute the sum of ranks for each group ($R_1$ and $R_2$).
    \item Compute the U statistics:
    \[
    U_1 = R_1 - \frac{n_1(n_1+1)}{2}, \quad U_2 = R_2 - \frac{n_2(n_2+1)}{2}
    \]
    \item Use the normal approximation for large samples:
    \[
    \mu_U = \frac{n_1 n_2}{2}, \quad \sigma_U = \sqrt{\frac{n_1 n_2 (n_1+n_2+1)}{12}}
    \]
    \[
    Z = \frac{U - \mu_U}{\sigma_U}
    \]
\end{enumerate}

\paragraph{Results.}

\begin{table}[H]
\centering
\caption{Mann--Whitney U Test Results (Independent Samples)}
\label{tab:mannwhitney}
\begin{tabular}{lcc}
\toprule
\textbf{$H_0$: $f_K = f_Q$ (or $f_\mathrm{Emergency Department} = f_\mathrm{Non-Emergency Department}$)} & \textbf{$p$-value} & \textbf{Decision} \\
\midrule
Length of Stay: Kings vs Queens           & $< 2.2 \times 10^{-16}$ & Reject $H_0$ \\
Total Charges: Kings vs Queens            & $< 2.2 \times 10^{-16}$ & Reject $H_0$ \\
Total Costs: Kings vs Queens              & $< 2.2 \times 10^{-16}$ & Reject $H_0$ \\
Birth Weight: Kings vs Queens             & $< 2.2 \times 10^{-16}$ & Reject $H_0$ \\
All Patient Refined Severity of Illness: Kings vs Queens  & $< 2.2 \times 10^{-16}$ & Reject $H_0$ \\
\midrule
\multicolumn{3}{l}{\textit{Additional comparison by Emergency Department Indicator:}} \\
Total Charges: Emergency vs Non-Emergency & $\approx 0$ & Reject $H_0$ \\
\bottomrule
\end{tabular}
\end{table}

\paragraph{Interpretation.}
Every single comparison came back highly significant:
\begin{itemize}[nosep]
    \item Total Charges between Kings and Queens had a p-value below $2.2 \times 10^{-16}$, essentially zero. There is no doubt that costs differ between the two boroughs.
    \item Total Costs, Length of Stay, Birth Weight, and All Patient Refined Severity of Illness were all significantly different between the counties as independent samples.
    \item When comparing Total Charges for Emergency vs.\ Non-Emergency admissions, that was also highly significant. Emergency visits clearly have a different cost profile.
\end{itemize}

\subsection{Test 5: Kolmogorov--Smirnov Two-Sample Test}

\paragraph{Test statistic.}
\[
  D = \sup_x |F_K(x) - F_Q(x)|,
  \qquad
  \sqrt{\tfrac{n_1 n_2}{n_1+n_2}}\, D \;\rightsquigarrow\; \text{Kolmogorov}.
\]

\begin{table}[H]
\centering
\caption{Kolmogorov-Smirnov Two-Sample Test: Kings vs Queens}
\label{tab:ks}
\begin{tabular}{lcr l}
\toprule
\textbf{Variable} & \textbf{D} & \textbf{$p$-value} & \textbf{Decision} \\
\midrule
Length of Stay & 0.0298 & $< 2.2\times 10^{-16}$ & Reject $H_0$ \\
Total Charges  & 0.0551 & $< 2.2\times 10^{-16}$ & Reject $H_0$ \\
Total Costs    & 0.1439 & $< 2.2\times 10^{-16}$ & Reject $H_0$ \\
Birth Weight   & 0.0464 & $< 2.2\times 10^{-16}$ & Reject $H_0$ \\
\bottomrule
\end{tabular}
\end{table}

\begin{table}[H]
\centering
\caption{Kolmogorov-Smirnov Two-Sample Test: Emergency vs Non-Emergency}
\label{tab:ks-ed}
\begin{tabular}{lcr l}
\toprule
\textbf{Variable} & \textbf{D} & \textbf{$p$-value} & \textbf{Decision} \\
\midrule
Length of Stay & 0.2959 & $< 2.2\times 10^{-16}$ & Reject $H_0$ \\
Total Charges  & 0.1950 & $< 2.2\times 10^{-16}$ & Reject $H_0$ \\
Total Costs    & 0.1674 & $< 2.2\times 10^{-16}$ & Reject $H_0$ \\
Birth Weight   & 0.0530 & $0.00155$ & Reject $H_0$ \\
\bottomrule
\end{tabular}
\end{table}

\paragraph{Interpretation.}
The entire cumulative distribution of every quantitative outcome
differs between the two counties and between Emergency and
Non-Emergency admissions, not merely their medians. Total Costs shows
the largest separation between Kings and Queens ($D=0.144$); Length of Stay shows
the largest separation between Emergency and Non-Emergency visits
($D=0.296$). These results subsume the Mann--Whitney findings, because
the Kolmogorov-Smirnov test is sensitive to differences in shape as well as
location---addressing the Fagerland and Sandvik (2009) critique that
the Mann--Whitney U alone can be misleading under unequal shapes.

\subsection{Test 6: Kruskal--Wallis Test}

\paragraph{Test statistic.}
\[
  H = \frac{12}{N(N+1)} \sum_{i=1}^{k} \frac{R_i^2}{n_i}
      \;-\; 3(N+1),
  \qquad H \sim \chi^2_{k-1} \text{ under } H_0.
\]

\begin{table}[H]
\centering
\caption{Kruskal--Wallis Test Results (all 15 testable outcome $\times$ grouping combinations)}
\label{tab:kw}
\begin{tabular}{l l r c r l}
\toprule
\textbf{Variable} & \textbf{Grouping} & \textbf{$\chi^2$} & \textbf{df} & \textbf{$p$-value} & \textbf{Decision} \\
\midrule
Length of Stay & Age Group        & 39{,}479 & 4 & $< 2.2\times 10^{-16}$ & Reject $H_0$ \\
Length of Stay & All Patient Refined Severity     & 82{,}402 & 4 & $< 2.2\times 10^{-16}$ & Reject $H_0$ \\
Length of Stay & Admission Type   & 25{,}198 & 5 & $< 2.2\times 10^{-16}$ & Reject $H_0$ \\
Length of Stay & Primary Payer    & 28{,}581 & 8 & $< 2.2\times 10^{-16}$ & Reject $H_0$ \\
Total Charges  & Age Group        & 65{,}223 & 4 & $< 2.2\times 10^{-16}$ & Reject $H_0$ \\
Total Charges  & All Patient Refined Severity     & 78{,}528 & 4 & $< 2.2\times 10^{-16}$ & Reject $H_0$ \\
Total Charges  & Admission Type   & 45{,}558 & 5 & $< 2.2\times 10^{-16}$ & Reject $H_0$ \\
Total Charges  & Primary Payer    & 35{,}087 & 8 & $< 2.2\times 10^{-16}$ & Reject $H_0$ \\
Total Costs    & Age Group        & 61{,}785 & 4 & $< 2.2\times 10^{-16}$ & Reject $H_0$ \\
Total Costs    & All Patient Refined Severity     & 76{,}947 & 4 & $< 2.2\times 10^{-16}$ & Reject $H_0$ \\
Total Costs    & Admission Type   & 50{,}005 & 5 & $< 2.2\times 10^{-16}$ & Reject $H_0$ \\
Total Costs    & Primary Payer    & 25{,}578 & 8 & $< 2.2\times 10^{-16}$ & Reject $H_0$ \\
Birth Weight   & All Patient Refined Severity     &  1{,}562 & 4 & $< 2.2\times 10^{-16}$ & Reject $H_0$ \\
Birth Weight   & Admission Type   &     40.5 & 4 & $< 2.2\times 10^{-16}$ & Reject $H_0$ \\
Birth Weight   & Primary Payer    &     52.6 & 7 & $< 2.2\times 10^{-16}$ & Reject $H_0$ \\
\bottomrule
\end{tabular}
\end{table}

\paragraph{Interpretation.}
Every outcome-by-grouping combination rejects $H_0$. Severity
consistently produces the largest $\chi^2$ statistic, indicating the
single strongest driver of distributional heterogeneity across Length of Stay,
Charges, and Costs---consistent with the findings of Stone
\textit{et~al.} (2022) and Nguyen \textit{et~al.} (2021).

\subsection{Test 7: Friedman Test}

\paragraph{Test statistic.}
\[
  Q = \frac{12}{nk(k+1)} \sum_{j=1}^{k} R_j^2 \;-\; 3n(k+1),
  \qquad Q \sim \chi^2_{k-1}\text{ under }H_0.
\]

With All Patient Refined Severity (4 levels) as treatment and each facility as a block,
only facilities observed at all four severity levels are retained.

\begin{table}[H]
\centering
\caption{Friedman Test: Severity as Treatment, Facility as Block}
\label{tab:friedman}
\begin{tabular}{l c c c c l}
\toprule
\textbf{Variable} & \textbf{\# Blocks} & \textbf{$k$} & \textbf{$\chi^2$} & \textbf{$p$-value} & \textbf{Decision} \\
\midrule
Length of Stay  & 5 & 4 & 13.29 & $0.0041$ & Reject $H_0$ \\
Total Charges   & 5 & 4 & 15.00 & $0.0018$ & Reject $H_0$ \\
Total Costs     & 5 & 4 & 15.00 & $0.0018$ & Reject $H_0$ \\
\bottomrule
\end{tabular}
\end{table}

\paragraph{Interpretation.}
Within any given hospital, Length of Stay, Charges, and Costs still differ
substantially across the four severity levels. The Friedman test
isolates the severity effect from facility-level differences in pricing
and case-mix, and confirms that illness severity is a fundamental
driver of inpatient utilisation. \textit{Note on block count}: only 5
of the facilities in the dataset had at least one discharge recorded
at all four All Patient Refined severity levels; this is a known limitation of the
Friedman design applied to administrative data, where low-volume
hospitals rarely see Extreme-severity cases. The test remains valid
and the results are conservative (fewer blocks reduce power), so
rejection of $H_0$ at $p<0.005$ is a robust finding despite the small
number of blocks.

\subsection{Test 8: Kendall's Tau}

\paragraph{Test statistic.}
\[
  \tau = \frac{C - D}{\binom{n}{2}},
  \qquad z = \frac{3\tau \sqrt{n(n-1)}}{\sqrt{2(2n+5)}}.
\]
To make the exact computation tractable on the full $N=368{,}495$
records, a random subsample of $n=5{,}000$ records was used.

\begin{table}[H]
\centering
\caption{Kendall's Tau on Six Pairs (subsample $n=5{,}000$)}
\label{tab:kendall}
\begin{tabular}{l r r l}
\toprule
\textbf{Pair} & \textbf{$\tau$} & \textbf{$p$-value} & \textbf{Decision} \\
\midrule
Length of Stay vs Charges     &  0.5761 & $< 2.2\times 10^{-16}$ & Reject $H_0$ \\
Length of Stay vs Costs       &  0.5858 & $< 2.2\times 10^{-16}$ & Reject $H_0$ \\
Length of Stay vs Birth Weight     & -0.1618 & $7.0\times 10^{-7}$    & Reject $H_0$ \\
Charges vs Costs   &  0.6708 & $< 2.2\times 10^{-16}$ & Reject $H_0$ \\
Charges vs Birth Weight & -0.1145 & $6.9\times 10^{-5}$    & Reject $H_0$ \\
Costs vs Birth Weight   & -0.1100 & $1.3\times 10^{-4}$    & Reject $H_0$ \\
\bottomrule
\end{tabular}
\end{table}

\subsection{Test 9: Spearman's Rank Correlation}

\paragraph{Test statistic.}
\[
  \rho_s = 1 - \frac{6\sum d_i^2}{n(n^2 - 1)}.
\]

\begin{table}[H]
\centering
\caption{Spearman's $\rho$ on Six Pairs (full data)}
\label{tab:spearman}
\begin{tabular}{l r r l}
\toprule
\textbf{Pair} & \textbf{$\rho_s$} & \textbf{$p$-value} & \textbf{Decision} \\
\midrule
Length of Stay vs Charges     &  0.7305 & $< 2.2\times 10^{-16}$ & Reject $H_0$ \\
Length of Stay vs Costs       &  0.7478 & $< 2.2\times 10^{-16}$ & Reject $H_0$ \\
Length of Stay vs Birth Weight     & -0.1733 & $< 2.2\times 10^{-16}$ & Reject $H_0$ \\
Charges vs Costs   &  0.8589 & $< 2.2\times 10^{-16}$ & Reject $H_0$ \\
Charges vs Birth Weight & -0.1257 & $< 2.2\times 10^{-16}$ & Reject $H_0$ \\
Costs vs Birth Weight   & -0.1490 & $< 2.2\times 10^{-16}$ & Reject $H_0$ \\
\bottomrule
\end{tabular}
\end{table}

\paragraph{Interpretation of Tests 8 and 9.} Charges and Costs are
strongly rank-associated ($\rho=0.86$). Length of Stay shows strong monotone
association with both Charges and Costs ($\rho\approx 0.73$--$0.75$).
Birth Weight shows only weak association with the other three
variables ($|\rho|<0.18$); the correlations are statistically
significant mainly because $N>300{,}000$. For every pair, Kendall's
$\tau$ is smaller in magnitude than Spearman's $\rho$, exactly as
predicted by Croux and Dehon (2010) under the bivariate-normal copula
approximation $\tau \approx (2/\pi)\arcsin(\rho)$.

\subsection{Test 10: Permutation Test (Monte Carlo, $B=5{,}000$)}

\paragraph{Test statistic.} For two samples $x_1,\dots,x_{n_1}$ and
$y_1,\dots,y_{n_2}$, using the median difference as the test
statistic,
\[
  T^{\text{obs}} = \tilde{X} - \tilde{Y},
\]
\[
  p = \frac{1}{B+1}\left(1 + \sum_{b=1}^{B}
    \mathbf{1}\bigl\{ |T^{(b)}| \ge |T^{\text{obs}}| \bigr\}\right).
\]

Using random subsamples of 3{,}000 records per county:

\begin{table}[H]
\centering
\caption{Permutation Test on Total Charges, Kings vs Queens ($B=5{,}000$)}
\label{tab:perm}
\begin{tabular}{l r r l}
\toprule
\textbf{Statistic} & \textbf{Observed} & \textbf{$p$-value} & \textbf{Decision} \\
\midrule
Median difference & \$2{,}858.07 & $0.011$  & Reject $H_0$ \\
\bottomrule
\end{tabular}
\end{table}

\paragraph{Interpretation.}
The permutation test is asymptotically equivalent to, but makes
strictly fewer assumptions than, the Mann--Whitney U test (Ernst, 2004;
Good, 2005). The median difference in Charges between Kings and
Queens is significant ($p=0.011$), so the null hypothesis of
exchangeable county labels is rejected. The robust median detects a
real distributional shift between the two counties even when extreme
outliers (the maximum single charge in Kings exceeds \$12 million)
are present in the sample. These findings reinforce the methodological
guidance of Chazard et al.\ (2017): on heavily right-skewed hospital
data, median-based or rank-based tests are more trustworthy summaries
of central location than methods that rely on the arithmetic average.

% ============================================================
% 6. COMPREHENSIVE SUMMARY (all 10 tests)
% ============================================================
\section{Comprehensive Summary}

Table~\ref{tab:summary-scope} gives a high-level view of what each
test covered and the overall decision at significance level $\alpha =
0.05$. The full eight-column row-per-test summary (objective question,
research question, null hypothesis, test statistic, test, p-value,
decision, and interpretation) has been moved to Section~4 and appears
there as Table~\ref{tab:objectives}.

\begin{table}[H]
\centering
\caption{Scope-Level Summary of All Ten Nonparametric Tests}
\label{tab:summary-scope}
\begin{tabular}{p{0.26\textwidth} p{0.32\textwidth} p{0.33\textwidth}}
\toprule
\textbf{Scope} & \textbf{Overall Decision} & \textbf{Test} \\
\midrule
7 numeric variables & Reject $H_0$ for every testable variable (Discharge Year not testable) & Runs Test \\
Length of Stay, Charges, Costs, Birth Weight, Severity (Kings vs Queens, paired) & Reject $H_0$ for all 5 variables & Sign Test \\
Length of Stay, Charges, Costs, Birth Weight, Severity (Kings vs Queens, paired) & Reject $H_0$ for all 5 variables & Wilcoxon Signed-Rank \\
Length of Stay, Charges, Costs, Birth Weight, Severity (Kings vs Queens) plus Charges (Emergency Department vs Non-Emergency Department) & Reject $H_0$ for all 6 comparisons & Mann-Whitney U \\
Kings vs Queens; Emergency Department vs Non-Emergency Department (4 variables each) & Reject $H_0$ for all 8 comparisons & Kolmogorov-Smirnov \\
4 outcomes $\times$ 4 groupings & Reject $H_0$ for all 15 testable combinations & Kruskal-Wallis \\
Length of Stay / Charges / Costs vs Severity & Reject $H_0$ for all 3 outcomes & Friedman (Facility block) \\
6 pairs of numeric variables & Reject $H_0$ for all 6 pairs & Kendall's $\tau$ \\
6 pairs of numeric variables & Reject $H_0$ for all 6 pairs & Spearman's $\rho$ \\
Kings vs Queens, median Charges & Reject $H_0$ for the median ($p=0.011$) & Permutation (Charges) \\
\bottomrule
\end{tabular}
\end{table}

% ============================================================
% 7. CONCLUSIONS
% ============================================================
\section{Conclusion, Comparison to Existing Work, and Future Research}

\subsection{Conclusions}

Ten nonparametric procedures were applied to Statewide Planning and Research Cooperative System 2022 discharge
records for Kings and Queens counties (368{,}495 admissions, 33
original variables plus 4 derived numeric columns). Nine of the ten
tests deliver a clean reject-$H_0$ verdict across essentially all
meaningful comparisons: the raw record ordering is non-random (Runs);
the median Length of Stay differs significantly between Kings and
Queens (Sign Test; $p<2.2\times 10^{-16}$); the paired distributions
of Length of Stay, Charges, and Costs differ between counties (Wilcoxon
Signed-Rank); the independent-sample distributions of Charges, Costs,
and Length of Stay differ between counties, and Charges differ between Emergency
and Non-Emergency admissions (Mann--Whitney U); the full cumulative
distributions of Length of Stay, Charges, Costs, and Birth Weight differ
significantly between the two counties and between Emergency and
Non-Emergency admissions (Kolmogorov-Smirnov); age, severity, admission type, and
payer all drive further heterogeneity (Kruskal--Wallis); severity
remains a dominant factor even after blocking by facility (Friedman);
and the four numeric outcomes are strongly rank-associated (Kendall,
Spearman), except that Birth Weight is only weakly related to the
other three. The Permutation Test confirms a significant median
Charges difference between counties ($p=0.011$), so the null
hypothesis of exchangeable county labels is rejected for the
distribution of Total Charges.

Overall, these results show that hospital utilization and costs really
do differ between Kings and Queens counties. There could be a lot of
reasons for this---different hospitals, different patient populations,
varying levels of disease severity, different insurance types, or
socioeconomic differences between the two boroughs. These findings
line up with what other researchers have found about geographic
variation in hospital outcomes within cities.

\subsection{Comparison to Existing Work}

The findings align with and extend the existing literature:

\begin{itemize}[nosep]
  \item Stone \emph{et al.}\ (2022) reported that Length of Stay for the same
    diagnosis can vary from 2 to 50+ days; the present
    Kruskal--Wallis result on Length of Stay across severity levels
    ($\chi^2=82{,}402$, $df=4$) quantifies this variability on a
    single-year, two-county slice of New York City.
  \item Chazard \emph{et~al.}\ (2017) found via simulation that
    Wilcoxon/permutation procedures have higher power than the
    $t$-test on skewed Length of Stay; the present permutation result
    on Charges bears this out directly---the median difference between
    counties is significant ($p=0.011$), and the rank-based
    Kolmogorov-Smirnov and Kruskal--Wallis tests on Length of Stay
    likewise reject on the full sample. This is a textbook
    demonstration of Chazard \emph{et~al.}'s recommendation to favour
    rank-based and median-based procedures on heavily right-skewed
    administrative data.
  \item Heslin \emph{et~al.}\ (2024) found significant weekday /
    weekend differences in Statewide Planning and Research Cooperative System admissions ($p<0.01$); the present
    Kolmogorov-Smirnov test on Emergency vs Non-Emergency ($D\ge 0.15$, $p<10^{-16}$)
    points to an even larger shape difference between these two
    admission modes.
  \item Croux and Dehon (2010) predicted $\tau<\rho$ under normal
    dependence; this is precisely reproduced in every pair analysed here.
  \item Qualls \emph{et~al.}\ (2010) argued that 43\% of Emergency Department-Length of Stay
    studies incorrectly use parametric methods; the present project,
    by using ten nonparametric tests, avoids this methodological
    pitfall entirely.
\end{itemize}

\subsection{Future Research and Innovative Extensions}

\begin{enumerate}[nosep]
  \item \textbf{All-borough extension}: generalise the analysis to
    Statewide Planning and Research Cooperative System 2023 covering all five New York City boroughs, permitting a
    six-group Kruskal--Wallis and a block-randomised Friedman with
    borough as treatment and facility as block.
  \item \textbf{Ordered-alternatives test}: apply the
    Jonckheere--Terpstra test to check the specific monotone trend of
    Length of Stay/Charges across the four ordinal severity levels.
  \item \textbf{Semi-parametric modelling}: follow Li \emph{et~al.}\
    (2022) to fit a semi-parametric time-to-discharge model,
    combining a parametric hazard with a nonparametric baseline.
  \item \textbf{Permutation-based regression}: build a nonparametric
    regression of Charges on facility, payer, severity, and
    demographics, with $p$-values derived from permuting the
    response; this extends the Permutation Test beyond a simple
    two-sample setting.
  \item \textbf{Bootstrap confidence intervals for effect sizes}:
    augment the $p$-value-based conclusions with non-parametric
    bootstrap 95\% confidence intervals (Efron \& Tibshirani, 1993) for key quantities
    such as Spearman's $\rho$ between Length of Stay and Charges and the median
    Charges difference between counties. Because $N>350{,}000$,
    standard $p$-values are essentially zero for any non-null effect;
    bootstrap confidence intervals would express the \emph{practical} size of each
    effect in its natural units (dollars, correlation), partially
    addressing the ``$p$-value inflation'' problem well known in very
    large administrative datasets.
  \item \textbf{Bootstrap quantile curves}: extend the bootstrap CI
    idea above to every quantile of the Charges distribution by
    county, producing a ``quantile shift'' curve that is fully
    nonparametric.
\end{enumerate}

% ============================================================
% 8. R CODE
% ============================================================
\section{R Code}

All analyses were executed in R~(4.3+). The complete code is included
in the accompanying \texttt{Merged\_Report.Rmd}. Key functions:

\begin{itemize}[nosep]
  \item custom \texttt{runs\_test()} using the normal approximation
    (\S\ref{sec:runs});
  \item \texttt{binom.test()} for the Sign Test;
  \item \texttt{wilcox.test(..., paired = TRUE)} for the Wilcoxon
    Signed-Rank Test;
  \item \texttt{wilcox.test(...)} for the Mann--Whitney U Test;
  \item \texttt{ks.test()} for the Kolmogorov--Smirnov Test;
  \item \texttt{kruskal.test()} for the Kruskal--Wallis Test;
  \item \texttt{friedman.test()} with \texttt{reshape()} wide-format
    construction for the Friedman Test;
  \item \texttt{cor.test(..., method = "kendall")} and
    \texttt{cor.test(..., method = "spearman")} for Kendall's Tau
    and Spearman's Rank Correlation;
  \item custom \texttt{perm\_two\_sample()} for the Permutation Test.
\end{itemize}

% ============================================================
% 9. REFERENCES (merged, deduplicated)
% ============================================================
\newpage
\section{References}

\begin{enumerate}[label={[\arabic*]}, leftmargin=*, nosep]

\item Anhoj, J., \& Olesen, A.~V. (2014). Run charts revisited:
  A simulation study of run chart rules for detection of non-random
  variation in health care processes. \textit{PLOS ONE}, 9(11),
  e113825. \url{https://doi.org/10.1371/journal.pone.0113825}

\item Bujang, M.~A., \& Sapri, F.~E. (2018). An application of the
  runs test to test for randomness of observations obtained from a
  clinical survey in an ordered population. \textit{Malaysian Journal
  of Medical Sciences}, 25(4), 146--151.
  \url{https://doi.org/10.21315/mjms2018.25.4.15}

\item Chazard, E., Ficheur, G., Beuscart, J.-B., \& Preda, C. (2017).
  How to compare the length of stay of two samples of inpatients?
  A simulation study to compare Type~I and Type~II errors of 12
  statistical tests. \textit{Value in Health}, 20(7), 992--998.
  \url{https://doi.org/10.1016/j.jval.2017.02.009}

\item Conover, W.~J. (1999). \textit{Practical Nonparametric
  Statistics} (3rd~ed.). John Wiley \& Sons. ISBN:~978-0471160687.

\item Croux, C., \& Dehon, C. (2010). Influence functions of the
  Spearman and Kendall correlation measures. \textit{Statistical
  Methods \& Applications}, 19, 497--515.
  \url{https://doi.org/10.1007/s10260-010-0142-z}

\item Dieleman, J.~L., Cao, J., Chapin, A., \emph{et al.}\ (2020).
  US health care spending by payer and health condition,
  1996--2016. \textit{JAMA}, 323(9), 863--884.
  \url{https://doi.org/10.1001/jama.2020.0734}

\item Efron, B., \& Tibshirani, R.~J. (1993). \textit{An Introduction
  to the Bootstrap}. Chapman \& Hall/CRC. ISBN:~978-0412042317.

\item Ernst, M.~D. (2004). Permutation methods: A basis for exact
  inference. \textit{Statistical Science}, 19(4), 676--685.
  \url{https://doi.org/10.1214/088342304000000396}

\item Fagerland, M.~W., \& Sandvik, L. (2009). The
  Wilcoxon--Mann--Whitney test under scrutiny. \textit{Statistics in
  Medicine}, 28(10), 1487--1497.
  \url{https://doi.org/10.1002/sim.3561}

\item Gibbons, J.~D., \& Chakraborti, S. (2011). \textit{Nonparametric
  Statistical Inference} (5th~ed.). CRC Press.
  ISBN:~978-1420077612.

\item Good, P.~I. (2005). \textit{Permutation, Parametric, and
  Bootstrap Tests of Hypotheses} (3rd~ed.). Springer.
  ISBN:~978-0387202792.

\item Heslin, S.~M., Henry, M., Litvak, E., Singer, A.~J., Thode, H.,
  \& Viccellio, P. (2024). An analysis of New York data:
  Fluctuations in hospital capacity are driven by variability in
  elective admissions and discharge activity. \textit{Cureus},
  16(4), e58404. \url{https://doi.org/10.7759/cureus.58404}

\item Hollander, M., Wolfe, D.~A., \& Chicken, E. (2013).
  \textit{Nonparametric Statistical Methods} (3rd~ed.). John Wiley
  \& Sons. ISBN:~978-0470387375.

\item Li, J., Couper, D., Hunt, K.~J., \& Mian, I.~U.~H. (2022).
  Semi-parametric time-to-event modelling of lengths of hospital
  stays. \textit{Journal of the Royal Statistical Society:
  Series~C (Applied Statistics)}, 71(5), 1667--1693.
  \url{https://doi.org/10.1111/rssc.12593}

\item Moore, B.~J., White, S., Washington, R., Coenen, N., \&
  Elixhauser, A. (2017). Identifying increased risk of readmission
  and in-hospital mortality using hospital administrative data:
  The AHRQ Elixhauser comorbidity index. \textit{Medical Care},
  55(7), 698--705.
  \url{https://doi.org/10.1097/MLR.0000000000000735}

\item Nguyen, D.~T., Ho, P.~D., Nguyen, T.~C., \& Van, N.~T.~C.
  (2021). Statistical analysis on length of stay in hospital.
  \textit{VNUHCM Journal of Engineering and Technology}, 3(SI3),
  SI97--SI106. \url{https://doi.org/10.32508/stdjet.v3iSI3.651}

\item Ning, X., Harkness, T., \& Gao, J. (2017). Spatio-temporal
  analysis for New York State SPARCS data. \textit{AMIA Annual
  Symposium Proceedings}, 2017, 1334--1343.
  \url{https://www.ncbi.nlm.nih.gov/pmc/articles/PMC5543354/}

\item Qualls, M., Pallin, D.~J., \& Schuur, J.~D. (2010). Parametric
  versus nonparametric statistical tests: The length of stay example.
  \textit{Academic Emergency Medicine}, 17(10), 1113--1121.
  \url{https://doi.org/10.1111/j.1553-2712.2010.00874.x}

\item Stone, K., Zwiggelaar, R., Jones, P., \& Mac~Parthal\'{a}in, N.
  (2022). A systematic review of the prediction of hospital length
  of stay: Towards a unified framework. \textit{PLOS Digital Health},
  1(4), e0000017.
  \url{https://doi.org/10.1371/journal.pdig.0000017}

\item Swed, F.~S., \& Eisenhart, C. (1943). Tables for testing
  randomness of grouping in a sequence of alternatives.
  \textit{Annals of Mathematical Statistics}, 14(1), 66--87.
  \url{https://doi.org/10.1214/aoms/1177731494}

\item Wilcoxon, F. (1945). Individual comparisons by ranking methods.
  \textit{Biometrics Bulletin}, 1(6), 80--83.
  \url{https://doi.org/10.2307/3001968}

\end{enumerate}


\end{document}
